%% How much of a hyperbolic orbifold does heat hear?
%% Springer Nature LaTeX template v3.1 (December 2024), sn-jnl class, numbered
%% Math and Physical Sciences reference style. No font package is loaded: the text and mathematics
%% are set in the template's default Computer Modern, which the figures match. The electronic
%% supplement is supplement.tex (its own PDF); references to it ("Table S5") are read from
%% build/supplement.aux by xr, so a clean build runs latexmk twice (BUILD.md).
\documentclass[pdflatex,sn-mathphys-num]{sn-jnl}

\usepackage{graphicx}%
\usepackage{amsmath,amssymb,amsfonts}%
\usepackage{amsthm}%
\usepackage{mathrsfs}%
\usepackage[title]{appendix}%
\usepackage{xcolor}%
\usepackage{booktabs}%
\usepackage{array}%
\usepackage{enumitem}%
\usepackage{xr}%
\externaldocument[S-]{build/supplement}% labels of the electronic supplement

\graphicspath{{../../}}% figures are addressed relative to the repository root

%% Theorem-like environments: one counter per section; Theorems A-C are lettered.
\theoremstyle{thmstyleone}%
\newtheorem{theorem}{Theorem}[section]%
\newtheorem{lemma}[theorem]{Lemma}%
\newtheorem{proposition}[theorem]{Proposition}%
\newtheorem{corollary}[theorem]{Corollary}%
\newtheorem{lettered}{Theorem}%
\renewcommand{\thelettered}{\Alph{lettered}}%
\theoremstyle{thmstylethree}%
\newtheorem{definition}[theorem]{Definition}%
\newtheorem{remark}[theorem]{Remark}%
\newtheorem{example}[theorem]{Example}%
\newtheorem{problem}{Problem}%

%% Notation (theory/CONVENTIONS.md).
\newcommand{\Orb}{\mathcal{O}}
\newcommand{\Cl}{\mathcal{C}}
\newcommand{\Sig}{\mathrm{Sig}}
\newcommand{\Kiso}{K_{\mathrm{iso}}}
\newcommand{\Kmult}{K_{\mathrm{mult}}}
\newcommand{\Area}{\operatorname{Area}}
\newcommand{\Teich}{\operatorname{Teich}}
\newcommand{\keytwo}{\operatorname{key}_2}
\newcommand{\Hyp}{\operatorname{Hyp}}
\newcommand{\cI}{\mathcal{I}}
\newcommand{\Hur}{\operatorname{Hur}}
\newcommand{\cond}{\operatorname{cond}}
\newcommand{\amp}{\operatorname{amp}}
\newcommand{\dist}{\operatorname{dist}}
\newcommand{\sepn}{\operatorname{sep}}
\newcommand{\gap}{\operatorname{gap}}
\newcommand{\diam}{\operatorname{diam}}
\newcommand{\Q}{\mathbb{Q}}
\newcommand{\R}{\mathbb{R}}
\newcommand{\C}{\mathbb{C}}
\newcommand{\Z}{\mathbb{Z}}
\newcommand{\cZ}{\mathrm{Z}}
\newcommand{\mult}{\operatorname{mult}}
\newcommand{\sref}[1]{\ref{S-#1}}% a label of the electronic supplement
\newcommand{\repo}[1]{\nolinkurl{#1}}% a path in the repository (breaks at slashes)
%% DOIs containing an underscore print correctly (CITATIONS.md, bibliography fixes).
\newcommand{\doiurl}[1]{\expandafter\url\expandafter{\detokenize{https://doi.org/#1}}}

%% Figures: each figure file is figures/out/F<k>.pdf in the repository, drawn at its final
%% width of 119 mm (figures/SPEC.md) and included at that size. The captions are the macros
%% \figcapOne ... \figcapNine of figures/captions.tex, attached to F1 ... F9 by \setfigcap.
\makeatletter
\newcommand{\setfigcap}[2]{\expandafter\def\csname figcap@#1\endcsname{#2}}
\newcommand{\figcap}[1]{\@ifundefined{figcap@#1}%
  {\@latex@error{No caption for figure #1 in figures/captions.tex}\@ehc}%
  {\csname figcap@#1\endcsname}}
\makeatother
% Captions of the figures, one \caption-ready macro per figure, each at most two short sentences:
% what is shown, and the one point to take from it; the decoding a reader needs (line styles,
% markers, scales) is in the sentence of the text that cites the figure. Notation: theory/CONVENTIONS.md.
% Colours are named by tone only ("dark", "light") and only where needed to read the figure, so the
% captions hold in greyscale print. Needs \usepackage{amssymb} (\mathfrak). The \ref labels are those
% of paper/jga/manuscript.tex, which reads this file and attaches \figcapOne ... \figcapEight to
% F1-F8 (\sref is its macro for a label of the electronic supplement). F9 belongs to the companion
% note, paper/arith/note.tex, which carries the same caption in its own notation. Every statement in
% a caption is asserted by the figure's script in figures/src/ or proved in the paper. Old and new
% word counts: paper/jga/FIGURE-RESTORE.md.

% F1 (Introduction) -- figures/out/F1.pdf
\newcommand{\figcapOne}{One triangle of $\mathcal O(2,8,8)$ (a) and of $\mathcal O(3,3,12)$ (b), in schematic shapes, coloured by $4\pi t\,\mathfrak h_t(x,x)$ at $t=0.02$. The two orbifolds share their first two heat invariants, yet heat lingers differently near their cone points.}

% F2 (Section 5) -- figures/out/F2.pdf
\newcommand{\figcapTwo}{Exact tilings of the hyperboloid by the triangle groups of $\mathcal O(2,8,8)$ (a) and $\mathcal O(3,3,12)$ (b), in an oblique orthographic view. The coloured triangle and its pale mirror image form one copy of each orbifold.}

% F3 (Section 3) -- figures/out/F3.pdf
\newcommand{\figcapThree}{The mirror argument of Theorem~\ref{thm:sigsep} for $(2,8,8)$, $(3,3,12)$ (a) and $(1;15)$, $(0;3,3,5,5)$ (b). Discs mark $X^*$ and rings $-X^*$; they coincide only for an orbifold against itself.}

% F4 (Section 3) -- figures/out/F4.pdf
\newcommand{\figcapFour}{The largest $K_{\mathrm{mult}}$ on complete equal-area classes and on constructed pairs, against $s=\mathrm{Area}/2\pi$. The open growth question lies between the two step curves.}

% F5 (Section 7) -- figures/out/F5.pdf
\newcommand{\figcapFive}{(a) The difference of the heat traces of $\mathcal O(2,8,8)$ and $\mathcal O(3,3,12)$ from computed spectra, against its cone-point term. (b) Differences of heat traces in the family $(0;3,3,3,3)$, against the predicted geodesic term.}

% F6 (Section 4, what heat does not hear) -- figures/out/F6.pdf
\newcommand{\figcapSix}{(a) Four members of the family $(0;3,3,3,3)$ with their shortest closed geodesics; (b) its eigenvalues below $120$ at eight moduli. All members share every heat invariant, yet the spectrum moves.}

% F7 (Section 5) -- figures/out/F7.pdf
\newcommand{\figcapSeven}{Reciprocal sums $R$ of the hyperbolic triads of sum $S$, one band for each least order $p=2,\dots,8$. Adjacent bands first overlap at the dots and first collide at the circles.}

% F8 (Section 6) -- figures/out/F8.pdf
\newcommand{\figcapEight}{Recovery error of the cone orders against a worst-case error $\delta$ in every heat invariant, for $(2,3,7)$, $(2,8,8)$ and $(4,4,4)$. The slopes are the exponents of Theorem~\ref{thm:S3}.}

% F9 (companion note, Section 6) -- figures/out/F9.pdf; the note's own copy uses \N and \ciso.
\newcommand{\figcapNine}{(a) The real locus of $C_{27/2}$ in the chart $x+y+z=1$, with the permutations of
$(1,4,4)$ and $(1,1,4)$, that is, $(2,8,8)$ and $(3,3,12)$ up to scale. (b) The number $\mathcal N(S)$ of pairs of sum at most $S$ sharing $c_1,c_2$, divided by $S$
(heavy), and its fit (dashed), far above the isosceles lower bound, which grows like
$(c_{\mathrm{iso}}+o(1))\log S$ (thin).}

% E1-E3: figures of the eigen manuscript (paper/eigen/manuscript.tex), same conventions.
% E1 (Section 8) -- figures/out/E1.pdf
\newcommand{\figcapEOne}{(a) The eigenvalues $\lambda_1,\dots,\lambda_6$ of $\mathcal O(2,3,m)$ against $m$, with their upper bounds; (b) the same eigenvalues, rescaled as $h_m^2(\lambda_j-\frac14)/\pi^2$. As the cone point of order $m$ becomes cusp-like, the low spectrum crowds towards $\frac14$.}

% E2 (Section 6) -- figures/out/E2.pdf
\newcommand{\figcapETwo}{The gaps between the signature terms of the heat trace of $\mathcal O(2,8,8)$ and of its five competitors, with the error of the geodesic term and the error of keeping finitely many eigenvalues. Where half the nearest gap exceeds both errors, finitely many eigenvalues decide the signature.}

% E3 (Section 7) -- figures/out/E3.pdf
\newcommand{\figcapEThree}{Eigenvalues needed to decide the signature on computed spectra, observed and certified, against the a-priori count of Theorem~\ref{thm:E}, by systole. The certificate needs $10^2$ to $10^{11}$ times fewer eigenvalues.}

\setfigcap{F3}{\figcapThree}\setfigcap{F4}{\figcapFour}
\setfigcap{F5}{\figcapFive}\setfigcap{F7}{\figcapSeven}\setfigcap{F8}{\figcapEight}
\newcommand{\paperfigure}[2]{%
  \begin{figure}[t]
  \centering
  \includegraphics[width=119mm]{figures/out/#1.pdf}%
  \caption{\figcap{#1}}\label{#2}
  \end{figure}}

\raggedbottom
\hypersetup{pdftitle={How much of a hyperbolic orbifold does heat hear?},
  pdfauthor={Palaash Gang, Akshaj Agadi, Arjun Veluri, Jerry Wang, Jeremy Barreto, Aarin Chouthaiwale},
  pdfkeywords={heat invariants, hyperbolic orbifolds, inverse spectral problems}}

\begin{document}

\title[How much of a hyperbolic orbifold does heat hear?]{How much of a hyperbolic orbifold does heat hear?}

\author*[1]{\fnm{Palaash} \sur{Gang}}\email{palaash.gang@indusschoolpune.com}
\author[2]{\fnm{Akshaj} \sur{Agadi}}\email{a0708@tamu.edu}
\author[1]{\fnm{Arjun} \sur{Veluri}}\email{arjun.v@indusschoolpune.com}
\author[3]{\fnm{Jerry} \sur{Wang}}\email{jerryiw@uci.edu}
\author[4]{\fnm{Jeremy} \sur{Barreto}}\email{jeremy.barreto@ltps.info}
\author[5]{\fnm{Aarin} \sur{Chouthaiwale}}\email{aarinchouthaiwale@gmail.com}

\affil*[1]{\orgname{Indus International School}, \orgaddress{\city{Pune}, \country{India}}}
\affil[2]{\orgname{Texas A\&M University}, \orgaddress{\city{College Station}, \state{Texas}, \country{USA}}}
\affil[3]{\orgname{University of California, Irvine}, \orgaddress{\city{Irvine}, \state{California}, \country{USA}}}
\affil[4]{\orgname{Lawrence High School}, \orgaddress{\city{Lawrenceville}, \state{New Jersey}, \country{USA}}}
\affil[5]{\orgname{DriveChange Learning and Resource Centre}, \orgaddress{\city{Pune}, \country{India}}}

%% ORCID iDs: the class's \orcid macro needs a logo file that the template does not ship,
%% so the iDs are printed as a title-page note through the class's note mechanism.
\presentaddresstxt{ORCID iDs}
\presentaddress{P.~Gang 0009-0006-7448-9488; A.~Agadi 0009-0009-1378-0556; A.~Veluri 0009-0001-5940-3272; J.~Wang 0009-0004-1814-275X; J.~Barreto 0009-0006-6894-7323; A.~Chouthaiwale 0009-0008-3878-7183.}

\abstract{For closed orientable hyperbolic $2$-orbifolds with cone points, we ask how many heat invariants determine the genus and the cone orders. The first $\lfloor\mathrm{Area}/\pi\rfloor+4$ always do, and no number independent of the area does. The number needed grows at least like the square root of the area; whether it grows linearly is equivalent to the open question whether the Prouhet--Tarry--Escott function satisfies $N(k)=O(k)$. For spheres with $n$ cone points, $n$ invariants suffice, and $n-1$ fail on integer orders for $n=3,4$. Among triangle orbifolds, two invariants determine every orbifold of cone-order sum at most $17$ and first fail at sum $18$, on $\mathcal O(2,8,8)$ and $\mathcal O(3,3,12)$, a pair isolated by a rank-$0$ elliptic curve; two suffice again at exactly $38$ sums between $19$ and $4800$, the first being $19$ and $21$--$25$, and three always suffice. For spheres with known $n$, recovering the orders from perturbed heat invariants is Lipschitz at simple orders and H\"older of exponent $1/k$ at $k$-fold orders, with constants depending on the true orders.}

\keywords{Heat invariants, Hyperbolic orbifolds, Inverse spectral problems, Cone singularities, Power sums, Prouhet--Tarry--Escott problem}

\pacs[MSC Classification]{58J53 (primary); 58J50, 35K08, 57R18, 30F35, 11D72, 11G05}

\maketitle

\section{Introduction}\label{sec:intro}

Put a unit of heat at a point of a closed surface and let it spread. For a short time $t$ it explores a disc of radius about $\sqrt t$ and learns only what lies in that disc. On a closed hyperbolic $2$-orbifold $\Orb$ a smooth point looks like every other smooth point, because the curvature is $-1$ everywhere, and a cone point of order $m$ looks like a disc folded $m$ times around its centre. The trace of the heat kernel adds up what all the points have learned; it has an asymptotic expansion
\begin{equation}\label{eq:heatseries}
  \cZ_\Orb(t)=\operatorname{tr}e^{-t\Delta}\ \sim\ \sum_{j\ge1}c_j(\Orb)\,t^{\,j-2}\qquad(t\downarrow0),
\end{equation}
and the coefficients $c_j(\Orb)$, the heat invariants, are spectral invariants \cite{donnelly1976,dggw2008}; see also the erratum \cite{dggw2017erratum}, which concerns only \cite[Thm~5.1]{dggw2008}. We ask a quantitative version of Kac's question \cite{kac1966}: how many heat invariants are needed to hear the cone points and the genus of $\Orb$, how fast that number must grow, and what heat cannot hear at all.

Write $\sigma(\Orb)=(g;m_1,\dots,m_n)$ for the \emph{signature} (genus and multiset of cone orders), $H_k(\Orb)=(c_1,\dots,c_k)$, $\Sig$ for the class of all closed orientable hyperbolic $2$-orbifolds whose singular points are cone points (orbifolds with mirrors are not locally orientable, so never isospectral to these \cite[Thm~4.7]{richardsonstanhope2020}; non-orientable ones are not considered) and $\Sig_0$ for its genus-$0$ part, and let $\Kmult(\Orb;\Cl)$ be the least $k$ such that $H_k$ separates $\Orb$ from every member of $\Cl$ with a different signature (Definition~\ref{def:K}). Let
\[
  f(A)=\max\{\Kmult(\Orb;\Sig):\ \Orb\in\Sig,\ \Area(\Orb)\le A\},
\]
and let $N(k)$ be the least size of a nontrivial solution of degree $k$ of the Prouhet--Tarry--Escott problem: two distinct $s$-multisets of integers with equal power sums of every exponent $1,\dots,k$ \cite[\S2, p.~6]{borweiningalls1994}.

\begin{theorem}[Hearing the signature]\label{thm:intro-signature}
\leavevmode
\begin{enumerate}[label=(\roman*)]
\item For every $\Orb\in\Sig$, the first $\lfloor\Area(\Orb)/\pi\rfloor+4$ heat invariants determine the genus and the cone-order multiset of $\Orb$ among all closed orientable hyperbolic $2$-orbifolds, and the number is read off from $c_1=\Area/4\pi$. No number independent of the area suffices: $\sup_\Orb\Kmult(\Orb;\Sig)=\sup_\Orb\Kmult(\Orb;\Sig_0)=\infty$.
\item For $A\ge8\pi$,
\[
  \Big\lfloor\sqrt{\tfrac13\big(\tfrac A{2\pi}-1\big)}\Big\rfloor+2\ \le\ f(A)\ \le\ \Big\lfloor\frac A\pi\Big\rfloor+4 .
\]
For $0<\alpha\le1$, $f(A)\ge cA^\alpha$ for all large $A$ and some $c>0$ if and only if $N(k)\le Ck^{1/\alpha}$ for all $k\ge1$ and some $C$. In particular $f$ grows linearly if and only if $N(k)=O(k)$.
\item Among spheres with $n$ cone points the first $n$ heat invariants determine the cone orders. The first $n-1$ do not: for $n=3,4$ they fail on integer orders, and for every $n\ge2$, as functions of real orders (the cone coefficients are rational in the order), they are not injective near any $n$ distinct positive reals; for integer orders and $n\ge5$ this is open (Problem~\ref{prob:sharp}).
\end{enumerate}
\end{theorem}

The mechanism is the same throughout. Heat sees a handle only through the area it uses up, and a cone point through a fixed function of its order; the $j$-th invariant adds one new odd power sum $\sum_im_i^{2j-3}$. Comparing two orbifolds becomes a moment problem for the orders of the first together with the negatives of the orders of the second: an odd-power Prouhet--Tarry--Escott system. Triangle orbifolds $\Orb(p,q,r)$ have no moduli, so for them the count can be made exact.

\begin{theorem}[The rigid case]\label{thm:intro-rigid}
Let $\Orb(p,q,r)$ be a hyperbolic triangle orbifold.
\begin{enumerate}[label=(\roman*)]
\item If $p+q+r\le17$, the first two heat invariants determine $\Orb(p,q,r)$ up to isometry among all hyperbolic triangle orbifolds.
\item The orbifolds $\Orb(2,8,8)$ and $\Orb(3,3,12)$ share their first two heat invariants. They are the only such pair with cone-order sum $18$, and no pair has a smaller sum.
\item The first three heat invariants always determine $\Orb(p,q,r)$.
\item For every integer $k\ge1$ the orbifolds $\Orb(2k,8k,8k)$ and $\Orb(3k,3k,12k)$ share their first two heat invariants with each other and with no third triangle orbifold.
\end{enumerate}
\end{theorem}

The supplement shows the two orbifolds of part~(ii) coloured by how much heat stays at each point. Finally, a measured heat trace is never exact.

\begin{theorem}[Stability]\label{thm:intro-stability}
Let $m$ be the cone orders of a hyperbolic sphere (genus $0$) with a known number $n$ of cone points, and $\tilde c_1,\dots,\tilde c_n$ approximations to its heat invariants $c_1,\dots,c_n$ with error at most $\delta$. For $\delta$ below an explicit bound, the recovery map of Section~\ref{sec:stability} returns, for each distinct order $a$ of multiplicity $k_a$, exactly $k_a$ numbers within $C_a\delta^{1/k_a}$ of $a$; the bound and the $C_a$ are explicit functions of $m$. The exponent $1/k_a$ cannot be improved for arbitrary data. For data that are the heat invariants of positive real orders it is $\frac12$, and sharp, at a double order (any $n$) and when all $n\ge3$ orders are equal; other configurations are open. For integer orders and $\delta$ below an explicit $\delta_{\rm thm}(m)$, rounding returns $m$ exactly.
\end{theorem}

Three limits of scope. The errors are errors in heat invariants, not in eigenvalues. The constants are evaluated at the true orders, so in practice the theorem is used a posteriori: recover a candidate, then check it (Remark~\ref{rem:stabscope}). And although Theorem~\ref{thm:intro-signature}(i) shows that the heat invariants determine $(g,n)$, we give no recovery procedure when $(g,n)$ is unknown.

\emph{What heat does not hear.} Every heat invariant is a function of the signature (classical; Proposition~\ref{prop:locality}), so off the triangle orbifolds no finite number of them determines $\Orb$ up to isometry (Corollary~\ref{cor:Kinf}). The shape enters the heat trace at order $t^{-1/2}e^{-\ell^2/4t}$, $\ell$ the systole: Theorem~\ref{thm:quantlocality} bounds this term with an explicit but qualitative constant, and Theorem~\ref{thm:sharp} shows that the exponent and the factor $t^{-1/2}$ are attained.

\emph{Roadmap.} Theorem~\ref{thm:intro-signature} is proved in Corollary~\ref{cor:sigarea} and Theorems~\ref{thm:signonuniform}, \ref{thm:growth}, \ref{thmA}, \ref{thmC}; Theorem~\ref{thm:intro-rigid} in Theorems~\ref{thm:three}, \ref{thm:threshold}, \ref{thm:minimal}, \ref{thm:isolation}; Theorem~\ref{thm:intro-stability} in Theorems~\ref{thm:S3}, \ref{thm:S4}, Proposition~\ref{prop:sharpexp} and Remark~\ref{rem:triple}. The lettered Theorems A--C are the algebraic core of Section~\ref{sec:signature}.

\subsection{Prior work and what is new}\label{sec:prior}

That the Laplace spectrum of a closed orientable hyperbolic $2$-orbifold determines its genus and its cone orders is due to Dryden and Strohmaier \cite[Thm~1.1, Prop.~3.3]{drydenstrohmaier2009}, after finiteness results of Stanhope \cite[Main Thms~1--2]{stanhope2005} and Dryden \cite[Thm~4.5]{dryden2004}; see also \cite[Thm~1]{doylerossetti2011}. U\c{c}ar recovered the cone orders from the whole sequence of heat invariants at any fixed nonzero curvature, with the curvature given \cite[Cors~4.21(iv), 4.23; method of Thm~3.40]{ucar2017}, and for $\chi\ge0$ a single coefficient determines the orbifold type: the invariant $c$, twelve times the $t^0$ coefficient, is a complete topological invariant there \cite[Thms~5.14--5.15]{dggw2008}. In dimension four, Abreu, Dryden, Freitas and Godinho recovered the weights of a weighted projective plane $\mathbb{CP}^2(N_1,N_2,N_3)$ with isolated singularities from the heat invariants $a_0,a_1,a_2$ of the Laplacian on functions and on $1$-forms \cite[Thm~1 and \S6.1]{adfg2008}, the closest precedent we know for finitely many heat invariants determining the orders of singular points. Theorem~\ref{thm:intro-signature} adds a count. To our knowledge, $\lfloor\Area/\pi\rfloor+4$ is the first explicit finite number of heat invariants shown to determine the genus and the cone-order multiset uniformly over all closed orientable hyperbolic $2$-orbifolds. Part~(i) also shows that the dependence on the area cannot be removed; explicit witnesses can be built from Prouhet's partition \cite[\S5.1, Thm~6]{alloucheshallit1999}, \cite{borweiningalls1994}, and Section~\ref{sec:pte} builds much smaller ones. A finite but likewise non-uniform count is known for Euclidean triangles and their Dirichlet eigenvalues \cite{changdeturck1989}. Part~(ii) ties the growth of the number of invariants needed, in both directions, to the Prouhet--Tarry--Escott problem; this relocates the growth question rather than settling it. All known bounds on $N(k)$ are quadratic, which gives the exponent $\frac12$ of (ii), and any exponent above $\frac12$ would settle in a strong form the open problem $N(k)=o(k^2)$ \cite[\S6, Problem~3]{borweiningalls1994}.

The mechanism of Theorem~\ref{thmA} is classical. By Newton's identities, two $n$-multisets of complex numbers have equal odd power sums $P_1,P_3,\dots,P_{2n-1}$ if and only if they agree after deleting pairs $\{a,-a\}$; this is the symmetric form of the Prouhet--Tarry--Escott problem \cite[Prop.~1 and \S3]{borweiningalls1994}. For positive reals, uniqueness of $n$ distinct numbers with $n$ prescribed power sums goes back to Steinig \cite{steinig1971}; see \cite[Lemma~3.2 and the remark following it]{laurens2023}, and \cite[Prop.~24]{msw2022} for positive integer exponents. Theorem~\ref{thmA} replaces the top odd power sum by the reciprocal sum $R$, which is the quantity the heat invariants supply; the condition $\prod_{i<j}(m_i+m_j)\ne0$ excludes pairs $\{a,-a\}$ in $m$. The linear system of Theorem~\ref{thmB} is the analogue, with $R$, of the odd-power-sum Newton identities of \cite[\S3]{korobovbugaevskaya2016}. What we add is the replacement of $P_{2n-1}$ by $R$, the closed form of $\det M$, whose zero set is the Orlando locus, and the sharpness statements of Theorem~\ref{thmC}.

Since $0<R<1$ for a hyperbolic triangle orbifold, the invariant $c$ of Dryden, Gordon, Greenwald and Webb encodes exactly the first two heat invariants. They observe that $c$ ``does not seem sufficiently strong to distinguish among'' the triangular pillows (their term for triangle orbifolds) with $\chi<0$, and note that the full spectrum determines the cone orders at curvature $-1$ \cite[Rem.~5.16]{dggw2008}. Theorem~\ref{thm:intro-rigid} makes this observation precise. The invariant $c$ separates the hyperbolic triangle orbifolds with $p+q+r\le17$ and first fails on $\Orb(2,8,8)$ and $\Orb(3,3,12)$, and the third heat invariant always separates them. To our knowledge no earlier result locates where a fixed finite number of heat invariants first fails on a family of hyperbolic orbifolds. On the length side, Philippe showed that triangle groups $(2,p,q)$, those with a right angle, are determined up to isometry by their length spectrum among such groups \cite[Thm~A]{philippe2008}, and later that any two triangle groups $(r,p,q)$ with the same length spectrum coincide \cite{philippe2010gd}, \cite[Thm~3.1]{philippe2010tsg}, a proof that compares the first two or three lengths counted with multiplicity; for Euclidean triangles the first three heat invariants suffice \cite[Thm~1]{griesermaronna2013}. The degeneracies are positive rational points on the cubics $C_\Lambda$, which Bremner, Guy and Nowakowski study for integer $\Lambda$ \cite[p.~117]{bgn1993}. Their remark that solutions come in reciprocal pairs $(x,y,z)$, $(1/x,1/y,1/z)$ (loc.\ cit.) does not use integrality, and at $\Lambda=\frac{27}2$, which they do not treat, the minimal pair is such a pair, $(1:4:4)$ and $(4:1:1)$, rescaled to a common sum. Its isolation is the statement that $C_{27/2}$ has rank $0$. The arithmetic of the two-coefficient degeneracies (their number, families and classes of every size) is the subject of the companion manuscript \cite{companion}; no proof here depends on it.

We know of no earlier quantitative stability estimate for recovering cone orders from heat invariants. For a sphere with a known number of cone points, the estimates of Section~\ref{sec:stability} combine the linear system of Theorem~\ref{thmB} with Ostrowski's root-perturbation bound \cite[Th\'eor\`eme~XXX]{ostrowski1940}. The loss of Lipschitz stability at a $k$-fold order, with H\"older exponent $1/k$, is the familiar behaviour of coalescing roots, and $1/k$ is sharp for general data. The rate is milder than in Prony-type moment problems with unknown amplitudes, where an $l$-node cluster costs the exponent $1/(2l-1)$ \cite{aby2015,bgy2020}, because here every weight equals $1$.

That heat invariants do not hear the shape is classical in substance. Selberg's trace formula for cocompact Fuchsian groups with elliptic elements \cite{hejhal1976}, in the form of \cite[eq.~(1)]{drydenstrohmaier2009} and applied to the heat kernel \cite{mckean1972,mckean1974corr}, \cite[Rem.~2.7]{garbinjorgenson2020}, writes the heat trace of a closed hyperbolic $2$-orbifold as three terms: an area term, one term for each cone point depending only on its order, and a sum over closed geodesics of size $O(t^{-1/2}e^{-\ell^2/4t})$. In particular every heat invariant is a function of the signature, as also follows from \cite[Thm~4.20]{ucar2017} at curvature $-1$. For surfaces the mechanism goes back to Huber and McKean \cite{huber1959,mckean1972}, and for hyperbolic orbisurfaces Dryden already read the systole off the rate $e^{-\ell^2/4t}$ \cite[proof of Thm~4.5]{dryden2004}. We record three consequences. No finite number of heat invariants determines an orbifold with moduli within its signature. The difference of two heat traces of one signature is bounded explicitly, with a constant depending on the area, the systole and the diameter. The prefactor $t^{-1/2}$ is attained when the systoles differ. Within one signature isospectral non-isometric pairs do exist \cite[Thm~A]{linowitzvoight2015}. The electronic supplement contains the full experiments and a blind recovery.

\section{Heat invariants of constant-curvature orbifolds}\label{sec:heat}

\subsection{Signatures and the trace formula}\label{sec:trace}

\begin{definition}[Signature]\label{def:signature}
Let $\Orb$ be a closed orientable $2$-orbifold whose singular points are cone points, of orders $m_1,\dots,m_n\ge2$, on an underlying surface of genus $g$. Its \emph{signature} is $\sigma(\Orb)=(g;m_1,\dots,m_n)$, the $m_i$ forming a multiset, and $\chi(\Orb)=2-2g-\sum_i(1-1/m_i)$. It carries a hyperbolic metric exactly when $\chi(\Orb)<0$, and then \cite[13.3.4--13.3.5]{thurston1980}, \cite[Thm~3.2]{drydenstrohmaier2009}
\begin{equation}\label{eq:area}
  \Area(\Orb)=-2\pi\chi(\Orb)=2\pi\Big(2g-2+\sum_{i=1}^n\big(1-\tfrac1{m_i}\big)\Big).
\end{equation}
A \emph{triangle orbifold} $\Orb(p,q,r)$, $p\le q\le r$, is the case $(g;n)=(0;3)$.
\end{definition}

Throughout, $R=\sum_i1/m_i$, $P_k=\sum_im_i^k$, $S_1=P_1$, and $e_k$ are the elementary symmetric functions of the orders ($e_0=1$); a \emph{hyperbolic triad} is a triple with $R<1$. As in \eqref{eq:heatseries}, $c_{l+2}$ is the coefficient of $t^l$. $\Sig_{0,n}\subset\Sig_0$ is the class of hyperbolic spheres with $n$ cone points.

\begin{definition}[$\Kiso$ and $\Kmult$]\label{def:K}
For a class $\Cl$ and $\Orb\in\Cl$, $\Kiso(\Orb;\Cl)$ and $\Kmult(\Orb;\Cl)$ are the least $k\ge1$ such that every $\Orb'\in\Cl$ with $H_k(\Orb')=H_k(\Orb)$ is isometric to $\Orb$, respectively has $\sigma(\Orb')=\sigma(\Orb)$ ($\min\emptyset=\infty$). So $\Kmult\le\Kiso$.
\end{definition}

Write $\Orb=\Gamma\backslash\mathbb H^2$ with $\Gamma\subset PSL(2,\R)$ cocompact, $h_t(r)=e^{-t(1/4+r^2)}$ and $g_t(u)=e^{-t/4}e^{-u^2/4t}/\sqrt{4\pi t}=\frac1{2\pi}\int h_t(r)e^{-iru}dr$, the normalization of \cite[eq.~(1)]{drydenstrohmaier2009}. For hyperbolic $\gamma\in\Gamma$ let $\ell(\gamma)$ be its translation length and $\gamma=\gamma_0^k$ with $\gamma_0$ primitive. A \emph{closed geodesic} is a hyperbolic conjugacy class (it may pass through cone points); the \emph{systole} $\ell$ is the least $\ell(\gamma)$ over hyperbolic $\gamma\in\Gamma$, and $\diam$ is the diameter of $\Orb$.

\begin{theorem}[Trace formula for the heat function]\label{thm:IEH}
For every $t>0$, $\cZ_\Orb(t)=\mathrm I(t)+\mathrm E(t)+\Hyp(t)$, all series and integrals converging absolutely, where $\mathrm E=\sum_i\mathrm E_{m_i}$ and
\begin{align*}
  \mathrm I(t)&=\frac{\Area(\Orb)}{4\pi}\int_\R r\tanh(\pi r)\,h_t(r)\,dr,\qquad
  \Hyp(t)=\sum_{[\gamma]\ \mathrm{hyperbolic}}\frac{\ell(\gamma_0)\,g_t(\ell(\gamma))}{2\sinh(\ell(\gamma)/2)}>0,\\
  \mathrm E_m(t)&=\sum_{j=1}^{m-1}\frac1{2m\sin(\pi j/m)}\int_\R\frac{e^{-2\pi jr/m}\,h_t(r)}{1+e^{-2\pi r}}\,dr .
\end{align*}
\end{theorem}

\begin{proof}
This is the Selberg trace formula for cocompact Fuchsian groups with elliptic elements as stated by Dryden and Strohmaier \cite[eq.~(1)]{drydenstrohmaier2009}, who take it from \cite{hejhal1976,iwaniec2002}; a cone point of order $m$ gives the elliptic classes $\rho^j$, $1\le j\le m-1$ \cite[p.~68]{drydenstrohmaier2009}. We use it for $h=\hat g$ with $g\in C_c^\infty(\R)$ even, the class \cite{drydenstrohmaier2009} intend. Garbin and Jorgenson state it for the heat function with the same normalization \cite[Rem.~2.7, (2.8)]{garbinjorgenson2020}; for completeness Lemma~\sref{lem:admissible} of the supplement gives the standard approximation argument, using Lemma~\ref{lem:counting}.
\end{proof}

\begin{lemma}[Counting closed geodesics]\label{lem:counting}
The number $n_\Orb(x)$ of closed geodesics of length at most $x$ satisfies $n_\Orb(x)\le2\pi(\cosh(x+3\diam)-1)/\Area(\Orb)\le\pi e^{x+3\diam}/\Area(\Orb)$.
\end{lemma}

\begin{proof}
Let $F$ be the Dirichlet domain at a point $x_0$ fixed by no nontrivial element; $F\subset\bar B(x_0,\diam)$. Every hyperbolic class has a representative $\gamma$ whose axis meets $F$, at $y$ say, and $d(x_0,\gamma x_0)\le2d(x_0,y)+d(y,\gamma y)\le\ell(\gamma)+2\diam$. The translates $\gamma F$ with $d(x_0,\gamma x_0)\le x+2\diam$ have disjoint interiors, area $\Area(\Orb)$ each, and lie in $B(x_0,x+3\diam)$.
\end{proof}

For closed surfaces Huber gave the asymptotic count \cite[Satz~9]{huber1959} and Buser a bound depending on the genus alone \cite{buser1992}.

\begin{lemma}[The hyperbolic term is small]\label{lem:hypbound}
For $0<t\le\ell^2/(2(1+\ell))$,
\[
  0<\Hyp(t)\le B(\ell,\diam,t):=\frac{\pi e^{3\diam}\,\ell e^{\ell/2}}{\Area(\Orb)\,(1-e^{-\ell})}\Big(1+\frac{2t}{\ell-t}\Big)\frac{e^{-\ell^2/4t}}{\sqrt{4\pi t}}\le\frac{C(\Area,\ell,\diam)}{t^{1/2}e^{\ell^2/4t}},
\]
\begin{equation}\label{eq:Cconst}
  C(A,\ell,\diam)=\frac{\sqrt\pi\,e^{3\diam}\,\ell\,e^{\ell/2}\,(2+3\ell)}{2A\,(1-e^{-\ell})\,(2+\ell)} .
\end{equation}
On this range $B$ increases with $\diam$ and decreases with $\ell$. In particular $\Hyp(t)=O(t^N)$ for every $N$.
\end{lemma}

\begin{proof}
With $\ell(\gamma_0)\le\ell(\gamma)$, $e^{-t/4}\le1$ and $2\sinh(x/2)\ge e^{x/2}(1-e^{-\ell})$ for $x\ge\ell$, $\Hyp(t)\le((1-e^{-\ell})\sqrt{4\pi t})^{-1}\int_{[\ell,\infty)}\phi\,dn_\Orb$, $\phi(x)=xe^{-x/2-x^2/4t}$. On the range, $(\log\phi)'=1/x-1/2-x/2t<0$ for $x\ge\ell$, so integrating by parts with Lemma~\ref{lem:counting}, $\int\phi\,dn_\Orb\le\frac{\pi e^{3\diam}}{\Area}\big(e^\ell\phi(\ell)+\int_\ell^\infty xe^{x/2-x^2/4t}dx\big)$, and since $x/(x-t)\le\ell/(\ell-t)$ on $[\ell,\infty)$ the integral is at most $\frac{2t\ell}{\ell-t}e^{\ell/2-\ell^2/4t}$. The logarithmic derivative of $B$ in $\ell$ is $-\frac12-(\frac\ell{2t}-\frac{1+\ell}\ell)-\frac{e^{-\ell}}{1-e^{-\ell}}-\frac{2t}{\ell^2-t^2}<0$, the bracket being nonnegative exactly on the range, which grows with $\ell$. Finally $1+2t/(\ell-t)\le(2+3\ell)/(2+\ell)$ on the range and $\pi/\sqrt{4\pi}=\sqrt\pi/2$.
\end{proof}

\subsection{The expansion for every order}\label{sec:heatexp}

The structure of the expansion is due to Donnelly and to Dryden, Gordon, Greenwald and Webb \cite{donnelly1976,dggw2008}, and U\c{c}ar computed every coefficient at constant curvature \cite{ucar2017}, starting from Watson's expansion for spherical lunes \cite{watson2005}, which he reproves with minor corrections; see also \cite[p.~2]{schueth2025}. We derive the coefficients again from Theorem~\ref{thm:IEH}, in closed form for every order. Write $\theta_j=\pi j/m$ and, for an integer $m\ge1$, $\Phi_m(u)=\sum_{j=1}^{m-1}(4m\sin\theta_j\sin(\theta_j-u))^{-1}$, analytic and even in $|u|<\pi/m$, with $\Phi_1=0$.

\begin{lemma}[Closed form]\label{lem:Phi}
For $0<|u|<\pi/m$, $\Phi_m(u)=\dfrac{\cot u-m\cot(mu)}{4m\sin u}$. Writing $u/\sin u=\sum_{i\ge0}\sigma_iu^{2i}$, so that $\sigma_0=1$ and $\sigma_i>0$, $\Phi_m(u)=\sum_k\phi_k(m)u^{2k}$ with
\begin{equation}\label{eq:phik}
  m\,\phi_k(m)=\frac14\sum_{n=1}^{k+1}\sigma_{k+1-n}\,\frac{4^n|B_{2n}|}{(2n)!}\,\big(m^{2n}-1\big),
\end{equation}
where $B_{2n}$ are the Bernoulli numbers.
\end{lemma}

\begin{proof}
$(\sin\theta\sin(\theta-u))^{-1}=(\cot(\theta-u)-\cot\theta)/\sin u$, and the $\cot\theta_j$ cancel in pairs $j,m-j$. The function $\sum_{j=0}^{m-1}\cot(x+\pi j/m)-m\cot(mx)$ is $\pi/m$-periodic, has no poles and tends to $0$ as $\operatorname{Im}x\to\pm\infty$, so it vanishes (Liouville); at $x=-u$ this gives the closed form. From $\frac x2\coth\frac x2=\sum_nB_{2n}x^{2n}/(2n)!$ at $x=2iu$, $\cot u-m\cot(mu)=\sum_{n\ge1}4^n|B_{2n}|(m^{2n}-1)u^{2n-1}/(2n)!$, using $\operatorname{sgn}B_{2n}=(-1)^{n+1}$ (Euler's formula for $\zeta(2n)$); multiply by $u/\sin u=\prod_n(1-u^2/n^2\pi^2)^{-1}$, which has positive coefficients, and by $1/(4mu)$.
\end{proof}

\begin{proposition}[The heat expansion at curvature $-1$]\label{prop:heatinput}
Let $\Orb\in\Sig$ have signature $(g;m_1,\dots,m_n)$. As $t\downarrow0$,
\[
  \cZ_\Orb(t)\sim\frac{\Area(\Orb)}{4\pi t}\sum_{k\ge0}\alpha_kt^k+\sum_{i=1}^n\sum_{l\ge0}b_l(m_i)\,t^l,
  \qquad \alpha_k=\frac{(-1)^k}{k!\,4^k}\sum_{l=0}^k\binom kl(-4)^lB_{2l}\big(\tfrac12\big),
\]
with Bernoulli polynomials $B_{2l}(x)$ and, for every integer $m\ge1$,
\begin{equation}\label{eq:bl}
  b_l(m)=\frac{(-1)^l\,p_l(m)}{m},\qquad
  p_l(m)=\frac1{4^l}\sum_{k=0}^{l}\frac{(2k)!}{k!\,(l-k)!}\;m\,\phi_k(m).
\end{equation}
Consequently $c_1=\Area(\Orb)/4\pi$ and, for $j\ge2$,
\begin{equation}\label{eq:cj}
  c_j(\Orb)=\alpha_{j-1}\,\frac{\Area(\Orb)}{4\pi}+\sum_{i=1}^nb_{j-2}(m_i).
\end{equation}
\end{proposition}

\begin{proof}
By Theorem~\ref{thm:IEH} and Lemma~\ref{lem:hypbound}, $\cZ_\Orb=\mathrm I+\mathrm E+O(t^N)$. For $0<a<2\pi$ let $F_a(r)=e^{-ar}/(1+e^{-2\pi r})$; Euler's beta integral gives $\int F_a(r)e^{sr}dr=[2\sin((a-s)/2)]^{-1}$ for $a-2\pi<s<a$, so the moments $\mathfrak m_n(a)=\int r^nF_a$ are its derivatives at $s=0$ (differentiation under the integral being dominated by $|r|^ne^{s_0|r|}F_a$ for $|s|\le s_0<\min(a,2\pi-a)$), and $|e^{-x}-\sum_{k<K}(-x)^k/k!|\le x^K/K!$ ($x=tr^2\ge0$) gives $\int F_ae^{-tr^2}dr\sim\sum_k\frac{(-t)^k}{k!}\mathfrak m_{2k}(a)$. Since $\mathrm E_m(t)=e^{-t/4}\sum_j(2m\sin\theta_j)^{-1}\int F_{2\theta_j}e^{-tr^2}dr$ and
\[
  \sum_{j=1}^{m-1}\frac{\mathfrak m_{2k}(2\theta_j)}{2m\sin\theta_j}=\frac{d^{2k}}{ds^{2k}}\Phi_m\Big(\frac s2\Big)\Big|_{s=0}=\frac{(2k)!}{4^k}\phi_k(m),
\]
multiplying by $e^{-t/4}=\sum_i(-t/4)^i/i!$ gives \eqref{eq:bl}. For $\mathrm I$, $r\tanh\pi r=|r|-2|r|/(e^{2\pi|r|}+1)$, $\int|r|e^{-tr^2}dr=1/t$, and the moments $\int_0^\infty r^{2k+1}(e^{2\pi r}+1)^{-1}dr=(1-2^{-2k-1})(-1)^kB_{2k+2}/(4(k+1))$ give $\alpha_k$, since $B_{2l}(\frac12)=(2^{1-2l}-1)B_{2l}$.
\end{proof}

\begin{lemma}[Cone polynomials]\label{lem:conepoly}
For every $l\ge0$, $p_l$ is an even polynomial with rational coefficients, of degree exactly $2l+2$, with $p_l(1)=0$, leading coefficient $|B_{2l+2}|/(2(l+1)!\,(2l+1))\ne0$, and $p_l(m)>0$ for every real $m>1$.
\end{lemma}

\begin{proof}
By \eqref{eq:phik} and \eqref{eq:bl}, $p_l$ is a combination with positive rational weights of the even polynomials $m^{2n}-1$, $1\le n\le l+1$; the degree $2l+2$ occurs only for $k=l$, $n=l+1$, with coefficient $\frac1{4^l}\cdot\frac{(2l)!}{l!}\cdot\frac{4^{l}|B_{2l+2}|}{(2l+2)!}=\frac{|B_{2l+2}|}{2(l+1)!(2l+1)}$.
\end{proof}

The first values are $\alpha_0,\dots,\alpha_4=1,-\frac13,\frac1{15},-\frac4{315},\frac1{315}$ and
\begin{equation}\label{eq:plexplicit}
  p_0(m)=\frac{m^2-1}{12},\qquad p_1(m)=\frac{m^4}{360}+\frac{m^2}{36}-\frac{11}{360},\qquad p_2(m)=\frac{m^6}{2520}+\frac{m^4}{720}+\frac{m^2}{180}-\frac{37}{5040}.
\end{equation}
At $K=-1$ the cone terms $b_0=p_0(m)/m$, $b_1=-p_1(m)/m$, $b_2=p_2(m)/m$ agree with Schueth's $a_0,a_1$ \cite[Rem.~4.2]{schueth2019} and $a_2$ \cite[Thm~4.1]{schueth2019}; she notes that the first two were already computed in \cite[\S5.6]{dggw2008} (there Prop.~5.5 and (5.10)). That $p_l(1)=0$ says that an ``order-$1$ cone point'' is a smooth point.

\begin{remark}[Agreement with U\c{c}ar; independence]\label{rem:ucaragree}\label{rem:proofC}
Putting $u=it/2$ in Lemma~\ref{lem:Phi} gives $m\Phi_m^{(2k)}(0)=2\cdot4^kk!\,m\,c^{\mathbb S}_k(\pi/m)$, with $c^{\mathbb S}_k$ the lune coefficient \cite[(4.25)]{ucar2017} (Watson's, as reproved there), so \eqref{eq:bl} is $(-1)^l$ times U\c{c}ar's cone contribution \cite[Thm~4.20(ii), with (4.33)--(4.34)]{ucar2017} at $\kappa=-1$, for every $l$; his smooth coefficients \cite[(4.35)]{ucar2017} are our $\alpha_k$ (also compared exactly for $l\le40$). The derivation is independent of the coefficient computations of \cite{dggw2008,ucar2017} and of the locality theorem of \cite{donnelly1976,dggw2008}: it uses only the trace formula for $\hat g$, $g\in C_c^\infty(\R)$ even, Lemma~\ref{lem:counting}, the approximation argument of Lemma~\sref{lem:admissible} of the supplement and classical analysis.
\end{remark}

\subsection{One odd power sum per coefficient}\label{sec:basis}

Put $\psi_k(x)=x^{2k-1}-x^{-1}$, $\Psi_k(m)=\sum_i\psi_k(m_i)=P_{2k-1}(m)-R(m)$ and $C_l(m)=\sum_ib_l(m_i)$.

\begin{lemma}[Triangular basis]\label{lem:triangular}
Write $p_l(x)=\sum_{k=0}^{l+1}a_{l,k}x^{2k}$. Then $p_l(x)/x=\sum_{k=1}^{l+1}a_{l,k}\psi_k(x)$ with $a_{l,l+1}\ne0$; hence for finite multisets $m,m'$ of positive reals, $C_l(m)=C_l(m')$ for $0\le l\le L-2$ if and only if $\Psi_k(m)=\Psi_k(m')$ for $1\le k\le L-1$.
\end{lemma}

\begin{proof}
$\sum_ka_{l,k}=p_l(1)=0$ gives the identity, and the system is triangular with diagonal $\pm a_{l,l+1}\ne0$ (Lemma~\ref{lem:conepoly}).
\end{proof}

So each power of the curvature switches on one more odd power sum: a cone point of order $m$ contributes at order $t^l$ a combination of $m^{2l+1},m^{2l-1},\dots,m,1/m$ with nonzero top term. For a sphere with $n$ cone points, $\Area/4\pi=(n-2-R)/2$, and the first $n$ heat invariants are an invertible triangular affine image of
\begin{equation}\label{eq:In}
  \cI_n(m)=(R,\ P_1,\ P_3,\ \dots,\ P_{2n-3}),
\end{equation}
$c_j$ being the first to involve $P_{2j-3}$ (Proposition~\ref{prop:frontend}). In particular $c_2=\chi/6+\sum_i(m_i^2-1)/(12m_i)$, the formula of \cite[Ex.~5.3, Prop.~5.5, (5.7)]{dggw2008}, and for a triangle orbifold
\begin{equation}\label{eq:s1inv}
  c_1=\frac{1-R}2,\qquad c_2=\frac{S_1+R-2}{12},\qquad c_3=\frac{1-R}{30}-\frac{P_3}{360}-\frac{S_1}{36}+\frac{11R}{360},
\end{equation}
so $c_1,c_2$ are equivalent to $(R,S_1)$ and $c_3$ adds $P_3$. A flat cone contributes nothing at order $t^1$: the cubic term is a curvature effect.

\begin{remark}[Other curvatures]\label{rem:curvature}
None of this is new. Proposition~\ref{prop:heatinput} holds at constant curvature $K$ with $K^l$ in place of $(-1)^l$ \cite[Thm~4.20]{ucar2017}; at $K=0$ the coefficients of $t^l$, $l\ge1$, vanish, and at $K=0,+1$ the $t^0$ coefficient determines the closed orientable orbifold type \cite[Thms~5.14--5.15, Prop.~5.22]{dggw2008}, as does, at $K=+1$, the full sequence of heat invariants \cite[Cor.~4.21(iv)]{ucar2017}. At $K=-1$ it no longer suffices (Theorem~\ref{thm:intro-rigid}(ii)).
\end{remark}

\section{Hearing the signature}\label{sec:signature}

Comparing two orbifolds means studying one signed multiset, the cone orders of the first together with the negatives of those of the second, whose reciprocal sum and low odd moments vanish. Everything here comes from one fact: such a multiset, if small compared with the number of vanishing moments, is symmetric under $x\mapsto-x$.

\subsection{The moment problem}

\begin{lemma}[Parity]\label{lem:parity}
In $\Q[x_1,\dots,x_N]$ let $s_j=\sum_ix_i^j$ and $e_j$ the elementary symmetric functions. For odd $j$, $e_j$ lies in the ideal generated by $s_1,s_3,\dots,s_j$, and $s_j$ in that generated by $e_1,e_3,\dots,e_j$.
\end{lemma}

\begin{proof}
By Newton's identities $e_j$ is a $\Q$-linear combination of products $s_{i_1}\cdots s_{i_r}$ with $\sum i_t=j$, and for odd $j$ some $i_t$ is odd; likewise for $s_j$.
\end{proof}

For a multiset $m$ of $n$ nonzero complex numbers let $\chi_m(z)=\prod_i(z+m_i)$ and $\Hur_k(\chi_m)$ its $k$-th Hurwitz determinant in the normalization of \cite[(1.37)]{holtztyaglov2012}; Orlando's formula \cite[Thm~1.17]{holtztyaglov2012} gives
\begin{equation}\label{eq:orlando}
  \Hur_{n-1}(\chi_m)=\prod_{i<j}(m_i+m_j).
\end{equation}

\begin{lettered}[Audibility of the cone orders]\label{thmA}
Let $m$ be a multiset of $n\ge1$ nonzero complex numbers with $m_i+m_j\ne0$ for all $i<j$, equivalently $e_n\Hur_{n-1}(\chi_m)\ne0$. If $m'$ is a multiset of $n$ nonzero complex numbers with $\cI_n(m')=\cI_n(m)$, then $m'=m$. In particular the first $n$ heat invariants determine the cone orders of a hyperbolic sphere with $n$ cone points among all such spheres: $\Kmult(\Orb;\Sig_{0,n})\le n$.
\end{lettered}

\begin{proof}
For $X=m\uplus(-m')$, $s_j(X)=P_j(m)+(-1)^jP_j(m')$ vanishes for odd $j\le2n-3$, so $e_j(X)=0$ for these $j$ (Lemma~\ref{lem:parity}); and $e_{2n-1}(X)=e_{2n}(X)(R(m)-R(m'))=0$ with $e_{2n}(X)=\pm e_n(m)e_n(m')\ne0$. So $\prod_{x\in X}(z-x)$ is even and $\mult_X(x)=\mult_X(-x)$, where $\mult_X(x)=\mult_m(x)+\mult_{m'}(-x)$ counts multiplicities. If $x$ lies in the support $A$ of $m$, then $-x\notin A$ by hypothesis, so $\mult_{m'}(x)=\mult_X(-x)=\mult_X(x)\ge\mult_m(x)$; summing over $A$, $n\le\sum_A\mult_{m'}\le n$, and $m'=m$. The equivalence is \eqref{eq:orlando}, and the heat invariants are an invertible affine image of $\cI_n$ (Section~\ref{sec:basis}).
\end{proof}

As in the classical case the hypothesis is sharp over $\C$: replacing a pair $a,-a$ in $m$ by $b,-b$ does not change $\cI_n(m)$.

\begin{lettered}[The linear system and its determinant]\label{thmB}
Let $T(z)=\tanh\big(\sum_{k\ \mathrm{odd}}P_kz^k/k\big)=\sum_kT_kz^k$. The equations
\begin{equation}\label{eq:linsys}
  e_{2j+1}-\sum_{i=0}^jT_{2i+1}\,e_{2j-2i}=0\quad(0\le j\le n-2),\qquad e_{n-1}-Re_n=0,
\end{equation}
with $e_k=0$ for $k>n$, form a square linear system $Me=b$ in $e=(e_1,\dots,e_n)$ whose coefficients are polynomials in $\cI_n$ and which the true $e$ solves. For $n\ge2$,
\[
  \det M=\varsigma_n\,\frac{\Hur_{n-1}(\chi_m)}{e_n}=\varsigma_n\,\frac{\prod_{i<j}(m_i+m_j)}{\prod_im_i},\qquad\varsigma_n=(-1)^{n(n+1)/2}.
\]
\end{lettered}

\begin{proof}
Let $f(z)=\prod_i(1+m_iz)=E(z^2)+zO(z^2)$. Since $\frac12\log\frac{f(z)}{f(-z)}=\sum_{k\ \mathrm{odd}}P_kz^k/k$, $T=\frac{f(z)-f(-z)}{f(z)+f(-z)}=\frac{zO(z^2)}{E(z^2)}$, and the coefficient of $z^{2j+1}$ in $zO=TE$ is \eqref{eq:linsys}; modulo $z^{2n-1}$ only $T_1,\dots,T_{2n-3}$, polynomials in $P_1,\dots,P_{2n-3}$, enter. For the determinant let $G:e\mapsto\cI_n$; differentiating $M(G(e))e=b(G(e))$ gives $M=-J\,\partial_eG$ with $J=\partial(Me-b)/\partial\cI_n$ at fixed $e$. In the column order $(R,P_1,\dots,P_{2n-3})$, $\partial T_k/\partial P_l=[z^{k-l}]\operatorname{sech}^2(U)/l$, $U=\sum_{k\ \mathrm{odd}}P_kz^k/k$, so after moving the last row $(-e_n,0,\dots,0)$ to the top $J$ is triangular with diagonal $-e_n$ and $-1/(2j+1)$. At distinct $m_i$, $\partial_eG=\partial_m\cI_n(\partial_me)^{-1}$; scaling column $i$ of $\partial_m\cI_n$ by $m_i^2$ gives rows $\varpi_rm_i^{2r}$ ($\varpi_0=-1$, $\varpi_r=2r-1$), a Vandermonde matrix in $m_i^2$, and $\det \partial_me=\prod_{i<j}(m_i-m_j)$. Collecting signs, $\det M=(-1)^{n+n(n-1)/2}\prod_{i<j}(m_i+m_j)/e_n$ on a dense set, hence wherever $e_n\ne0$.
\end{proof}

So recovering the orders is a linear problem followed by root-finding, singular exactly where Theorem~\ref{thmA} fails.

\begin{lettered}[The fibres of $n-1$ invariants]\label{thmC}
\leavevmode
\begin{enumerate}[label=(\arabic*)]
\item \emph{Pair criterion} (an odd symmetric system in the sense of \cite[\S3]{borweiningalls1994}). For multisets $m,m'$ of $n\ge2$ nonzero complex numbers and $Q(z)=\prod_i(z-m_i)\prod_j(z+m'_j)$, $\cI_{n-1}(m)=\cI_{n-1}(m')$ if and only if $Q(z)-Q(-z)=2\varkappa z^3$ for some $\varkappa\in\C$; if $m_i+m_j\ne0$ for $i<j$, then $\varkappa=0$ exactly when $m'=m$.
\item \emph{Real sharpness.} For $n\ge2$, every multiset of $n$ distinct positive reals lies on a real curve of multisets of positive reals sharing $\cI_{n-1}$.
\item \emph{Integer sharpness for $n=3,4$.} $\{2,8,8\}$ and $\{3,3,12\}$ share $\cI_2$ ($R=\frac34$, $P_1=18$) and have $P_3=1032\ne1782$; $\{3,10,15,30\}$ and $\{4,5,21,28\}$ share $\cI_3$ ($R=\frac8{15}$, $P_1=58$, $P_3=31402$) and have $P_5=25159618\ne21298618$. So $\Kmult(\Orb;\Sig_{0,n})=n$ for these orbifolds.
\end{enumerate}
\end{lettered}

\begin{proof}
(1) As in Theorem~\ref{thmA}, agreement gives $e_j(X)=0$ for odd $j\ne2n-3$, so $Q(z)-Q(-z)=-2e_{2n-3}(X)z^3$, and conversely by the second half of Lemma~\ref{lem:parity}; if $\varkappa=0$, $Q$ is even and $m'=m$. (2) Scaling column $i$ of the Jacobian of $\cI_{n-1}$ by $m_i^2$ gives a Vandermonde-type matrix of rank $n-1$ at distinct $m_i$, so the fibre is a smooth curve near $m$, whose other points are multisets different from $m$. (3) Direct computation; the four multisets are hyperbolic signatures, and spheres with $n$ cone points share $H_{n-1}$ exactly when they share $\cI_{n-1}$.
\end{proof}

\begin{remark}[Integer witnesses]\label{rem:witnesses}
Rational witnesses scale to hyperbolic integer ones, and a witness is disjoint, since a common order could be cancelled, contradicting Theorem~\ref{thmA}. For $n=4$, call a pair $(A,B)$ of $4$-multisets of nonzero integers a \emph{pencil splitting} if $e_1(A)=e_1(B)=0$, $e_3(A)=e_3(B)$, $e_4(A)=e_4(B)$ and $A\uplus(-B)$ has four entries of each sign and no pair $\{z,-z\}$. Then $s_1$, $s_3=3e_3$ and $s_{-1}=e_3/e_4$ agree, so the positive part of $A\uplus(-B)$ and the negative of its negative part form a witness ($A=\{-30,-3,5,28\}$, $B=\{-21,-4,10,15\}$ give that of Theorem~\ref{thmC}(3)). Pencil splittings with all entries of absolute value at most $130$ (respectively $220$) give exactly $17$ (respectively $49$) witnesses, of which $14$ (respectively $30$) are primitive. Not every witness arises so: of the $107$ primitive witnesses with all orders at most $440$, $6$ have no pencil splitting, the smallest being $(0;16,16,74,74)$ and $(0;11,37,44,88)$. For $n=5$ no witness has all orders at most $120$. The searches, their programs and the witnesses are in Section~\sref{sec:searches} of the supplement.
\end{remark}

\subsection{Heat data of a signature}

\begin{lemma}[Heat data of a signature]\label{lem:sigdata}
Let $\Orb,\Orb'\in\Sig$ have signatures $(g;m)$ and $(g';m')$ with $n$ and $n'$ cone points, and $L\ge1$. Then $H_L(\Orb)=H_L(\Orb')$ if and only if
\begin{equation}\label{eq:red1}
  2g+n-R(m)=2g'+n'-R(m')\quad\text{and}\quad\Psi_k(m)=\Psi_k(m')\ \ (1\le k\le L-1).
\end{equation}
Then $d=R(m')-R(m)=2(g'-g)+n'-n\in\Z$, and the paddings $U=m\uplus\{1\}^{\max(d,0)}$, $V=m'\uplus\{1\}^{\max(-d,0)}$ satisfy
\begin{equation}\label{eq:red2}
  R(U)=R(V),\qquad P_j(U)=P_j(V)\ \ (j\ \text{odd},\ j\le2L-3),\qquad |V|-|U|=2(g-g'),
\end{equation}
and $|U|+|V|=2\max(n+g-g',\,n'+g'-g)$. Conversely, if paddings by $1$s satisfy \eqref{eq:red2}, the two orbifolds have equal area and equal $H_L$.
\end{lemma}

\begin{proof}
$c_1$ agrees exactly when the areas do, the first equation; then $c_2,\dots,c_L$ agree exactly when $C_0,\dots,C_{L-2}$ do (Lemma~\ref{lem:triangular}). An order $1$ changes neither $n-R$ nor any $\Psi_k$, and the rest is direct computation, with $\Area(g;m)=2\pi(2g-2+|U|-R(U))$.
\end{proof}

\begin{theorem}[Separation]\label{thm:sigsep}
Let $\Orb,\Orb'\in\Sig$ have different signatures and equal $H_L$, and let $U^*,V^*$ be the paddings of Lemma~\ref{lem:sigdata} with their common elements removed (they do not depend on the paddings). Then $|U^*|+|V^*|\ge2L+2$. In particular, if $L\ge\max(n+g-g',n'+g'-g)$, equal $H_L$ implies equal signatures.
\end{theorem}

\begin{proof}
Admissible paddings differ by equally many $1$s on both sides. Cancelling common elements preserves \eqref{eq:red2}, so $U^*,V^*$ are disjoint, and $T=|U^*|+|V^*|\equiv2(g-g')$ is even. If $T\le2L$, the odd power sums of $X^*=U^*\uplus(-V^*)$ vanish up to $2L-3\ge T-3$ and $e_{T-1}(X^*)=e_T(X^*)(R(U^*)-R(V^*))=0$, so $\prod_{x\in X^*}(z-x)$ is even and $\mult_{U^*}=\mult_{V^*}$ on $(0,\infty)$; disjointness forces $U^*=V^*=\emptyset$, so the signatures agree. Hence $T\ge2L+1$, so $T\ge2L+2$; and $T\le|U|+|V|$.
\end{proof}

Figure~\ref{fig:F3} draws $X^*$ and its mirror image for two pairs sharing exactly $c_1,c_2$.

\paperfigure{F3}{fig:F3}

\begin{corollary}[Genus and cone orders from $\lfloor\Area/\pi\rfloor+4$ heat invariants]\label{cor:sigarea}
For every $\Orb\in\Sig$, $\Kmult(\Orb;\Sig)\le\lfloor\Area(\Orb)/\pi\rfloor+4$. A hyperbolic sphere of area $A$ has at most $A/\pi+4$ cone points, with equality exactly when all orders are $2$, so the cone count is determined without being known in advance.
\end{corollary}

\begin{proof}
Each cone point adds at least $\frac12$ to $\sum_i(1-1/m_i)$, so $n+4g\le\Area/\pi+4$, with equality in genus $0$ exactly when all orders are $2$. The same holds for $\Orb'$ of the same area, so $n+g-g'\le n+4g\le\Area/\pi+4$, likewise for $n'+g'-g$, and Theorem~\ref{thm:sigsep} applies.
\end{proof}

\begin{example}\label{ex:siggenus}
$(1;15)$ and $(0;3,3,5,5)$ share exactly two heat invariants: both areas are $2\pi\cdot\frac{14}{15}$ and $\Psi_1=\frac{224}{15}$ on both sides. Here $|U^*|+|V^*|=6=2L+2$, so Theorem~\ref{thm:sigsep} is attained; a handle costs no more invariants than a cone point.
\end{example}

\subsection{How many invariants: the Prouhet--Tarry--Escott connection}\label{sec:pte}

\begin{definition}[Configurations]\label{def:config}
An \emph{$L$-configuration} is a nonempty finite multiset $Z$ of nonzero rationals of even size, with no pair $\{z,-z\}$, such that $\sum_{z\in Z}z^j=0$ for odd $j\le2L-3$ and $\sum_{z\in Z}z^{-1}=0$. Its \emph{imbalance} is $\iota(Z)=\#\{z>0\}-\#\{z<0\}$.
\end{definition}

By Lemma~\ref{lem:sigdata} and Theorem~\ref{thm:sigsep}, two orbifolds in $\Sig$ with different signatures share their first $L$ heat invariants exactly when $Z=U^*\uplus(-V^*)$ is an $L$-configuration with imbalance $\iota(Z)=2(g'-g)$. Conversely, an $L$-configuration $Z$ scaled to coprime integers gives, with $U=Z_{>0}$, $V=-Z_{<0}$ and all entries $1$ removed, signatures $(g;U)$ and $(g';V)$ with $g'-g=\iota(Z)/2$ and equal area, the smaller genus being the least that makes its side hyperbolic (the other side then has the same positive area), sharing exactly $L$ heat invariants if and only if $\sum_Zz^{2L-1}\ne0$. Since $Z$ has no pair $\{z,-z\}$, $Z$ and $-Z$ form a solution of degree $2L-2$ of the Prouhet--Tarry--Escott problem, so $|Z|\ge N(2L-2)$.

\begin{theorem}[Descartes bound]\label{thm:descartes}
Every $L$-configuration has $|\iota(Z)|\le|Z|-2L$. Hence if orbifolds of genera $g\ne g'$ share their first $L$ heat invariants, then $|U^*|+|V^*|\ge2L+2|g-g'|$, and if moreover $g\ne g'$ and $|U^*|+|V^*|=2L+2$, the least size allowed by Theorem~\ref{thm:sigsep}, then $|g-g'|=1$ and $\{|U^*|,|V^*|\}=\{L,L+2\}$.
\end{theorem}

\begin{proof}
Let $T=|Z|$ and $Q(x)=\prod_{z\in Z}(x-z)=\sum_k(-1)^ke_kx^{T-k}$. By Lemma~\ref{lem:parity}, $e_k=0$ for odd $k\le2L-3$, and $e_{T-1}=e_T\sum_zz^{-1}=0$, so at most $(T-2L)/2$ coefficients of odd degree are nonzero. Let $V^\pm$ be the numbers of sign changes of $Q(x)$ and $Q(-x)$. Consecutive nonzero coefficients of degrees $d>d'$ contribute to $V^++V^-$ at most $d-d'$, and to $V^+-V^-$ zero if $d-d'$ is even and $\pm1$ if it is odd; so $V^++V^-\le T$. All roots are real and nonzero, so by Descartes' rule $\#\{z>0\}\le V^+$, $\#\{z<0\}\le V^-$ with total $T$; both are equalities and $\iota=V^+-V^-$. Each odd gap ends at a coefficient of odd degree, and each such coefficient ends at most two gaps, so $|\iota|\le T-2L$. Finally $\iota=2(g'-g)$.
\end{proof}

\begin{proposition}[Doubling]\label{prop:doubling}
Let $X\ne Y$ be $n$-multisets of positive integers with $P_j(X)=P_j(Y)$ for odd $j\le2L-3$, $U=X\uplus2Y\uplus2Y$, $V=Y\uplus2X\uplus2X$, and $\mult_X,\mult_Y$ the multiplicities of $1$ in $X,Y$. Then $(0;U\setminus\{1\})$ and $(0;V\setminus\{1\})$, every $1$ removed, are hyperbolic, with distinct signatures, equal area $<2\pi(3n-2)$, equal first $L$ heat invariants, and $3n-\mult_X$ and $3n-\mult_Y$ cone points.
\end{proposition}

\begin{proof}
For every $j$, including $j=-1$, $P_j(U)-P_j(V)=(1-2^{j+1})(P_j(X)-P_j(Y))$, and $|U|=|V|=3n$; so Lemma~\ref{lem:sigdata}, with $U,V$ as paddings, gives equal area $2\pi(|U|-R(U)-2)<2\pi(3n-2)$ and equal $H_L$, once both are hyperbolic. Here $n\ge2$; $V$ contains the $2n$ entries of $2X\uplus2X$, all $\ge2$, so $\sum(1-1/v)\ge n>2$ for $n\ge3$, and for $n=2$, $Y$ has at most one $1$, so $V$ has five entries $\ge2$; likewise for $U$. If $\nu(x)$ is the multiplicity of $x$ in $X$ minus that in $Y$ and $x^*$ the largest $x$ with $\nu(x)\ne0$, then $2x^*$ has multiplicity $-2\nu(x^*)\ne0$ in $U$ minus $V$.
\end{proof}

\begin{theorem}[No uniform number of heat invariants]\label{thm:signonuniform}
For every $L\ge2$ there are orbifolds in $\Sig$ of different genera, the smaller of which can be any integer $g\ge1$, and orbifolds in $\Sig_0$ with different numbers of cone points, sharing their first $L$ heat invariants. Hence $\sup_\Orb\Kmult(\Orb;\Sig)=\sup_\Orb\Kmult(\Orb;\Sig_0)=\infty$.
\end{theorem}

\begin{proof}
The genus pairs are built in the proof of Theorem~\ref{thm:growth}(e); adding equally many handles to both keeps the areas, hence the invariants, equal. For the cone count, take a nontrivial solution of degree $2L-3$ of the Prouhet--Tarry--Escott problem \cite[Prop.~3]{borweiningalls1994}, cancel common elements and translate so that the least element is $1$, in $X$ say; Proposition~\ref{prop:doubling} with $\mult_X\ge1$, $\mult_Y=0$ applies.
\end{proof}

For the growth, let $f_g(A)$ and $f_n(A)$ be the maxima over orbifolds of area at most $A$ of the least $k$ such that every $\Orb'\in\Sig$ (respectively every genus-$0$ $\Orb'$) with $H_k(\Orb')=H_k(\Orb)$ has the same genus (respectively the same number of cone points), $f_n$ being maximised over genus-$0$ $\Orb$; $f_g,f_n\le f$. One has $k+1\le N(k)\le\frac12k(k+1)+1$ \cite[Props.~2--3]{borweiningalls1994}; Melzak credits the pigeonhole bound on the right to Hardy and Wright \cite[p.~233]{melzak1961} and reports Wright's improvement $N(k)\le(k^2+4)/2$ (1935) \cite[p.~234]{melzak1961}, together with an exact but non-constructive formula and numerical bounds for $k\le29$ \cite[Table~1]{melzak1961}. None of these is $o(k^2)$, and $N(k)=o(k^2)$ remains open \cite[\S6, Problem~3]{borweiningalls1994}; Problem~4 there, $O(k^2)$ for the exact-degree analogue, was settled by Wooley; his bound for it is now $\frac12k(k+1)+1$ \cite[Thm~13.1]{wooley2019}. Ideal solutions, $N(k)=k+1$, are known only for $k\le9$ and $k=11$, and whether $N(k)=k+1$ for all $k$ is open \cite{cmsv2024}, \cite[\S1]{crootmaoyip2026}.

\begin{theorem}[Growth]\label{thm:growth}
\leavevmode
\begin{enumerate}[label=(\alph*)]
\item For every $L\ge2$ two hyperbolic genus-$0$ orbifolds with different signatures and area less than $2\pi(3(L-1)^2+1)$ share their first $L$ heat invariants. Consequently, for $A\ge8\pi$,
\[
  \Big\lfloor\sqrt{\tfrac13\big(\tfrac A{2\pi}-1\big)}\Big\rfloor+2\ \le\ f(A)\ \le\ \Big\lfloor\frac A\pi\Big\rfloor+4 .
\]
\item If $L\ge2$ and $f(A)\ge L+1$, then $N(2L-2)\le2\lfloor A/\pi\rfloor+8$. (c) If $\beta>0$ and $N(k)\le Ck^\beta$ for all $k\ge1$, then $f(A)\ge\frac12(A/6\pi C)^{1/\beta}$ for all $A\ge8\pi$.
\item[(d)] For $0<\alpha\le1$, $f(A)\ge cA^\alpha$ for all large $A$ and some $c>0$ if and only if $N(k)\le Ck^{1/\alpha}$ for all $k\ge1$ and some $C$. In particular $f(A)=\Theta(A)$ if and only if $N(k)=O(k)$.
\item[(e)] $f_g(A)\ge\sqrt{A/16\pi}$ for $A\ge16\pi$, $f_n(A)\ge\sqrt{A/12\pi}$ for $A\ge8\pi$, and (d) holds verbatim for $f_g$ and $f_n$.
\end{enumerate}
\end{theorem}

The proofs are in Appendix~\ref{app:pte}. With the known quadratic bounds on $N(k)$, (c) gives exactly the exponent $\frac12$ of (a), and any exponent above $\frac12$ in (d) would prove $N(k)=O(k^{2-\epsilon})$. Figure~\ref{fig:F4} compares the bounds with the exact largest $\Kmult$ on complete classes of signatures of equal area (its caption says which areas) and with the pairs below.

\paperfigure{F4}{fig:F4}

\begin{example}[Small pairs]\label{ex:ptepairs}
These pairs share exactly $L$ heat invariants (exact computation); all of them are written out in Table~\sref{tab:pairs} of the supplement. (i) $L=3$: $(0;3,10,15,30)$ and $(0;4,5,21,28)$, of area $2\pi\cdot\frac{22}{15}$. This is the $n=4$ witness of Theorem~\ref{thmC}(3). Its area is less than that of the smallest genus-changing pair sharing three invariants found here, $(1;15,15,15)$ and $(0;3,3,5,7,7,21)$, of area $2\pi\cdot\frac{14}5$. (ii) $L=3$, genus $0$ with $7$ and $8$ cone points: $(0;4,4,5,5,6,12,12)$ and $(0;2,2,2,3,10,10,10,10)$, the doubling of $[1,5,5]=[2,3,6]$ \cite[A.1.6]{chen2025survey} after cancelling the common order $6$. (iii) $L=4,\dots,7$: genus-$0$ pairs with equal cone counts and area less than $2\pi\cdot5,7,10,18$, from odd ideal symmetric solutions of size $7$ \cite[p.~9]{borweiningalls1994}, Letac's two solutions of size $9$ \cite[p.~2]{cmsv2024}, equal sums of odd powers \cite[A.1.33]{chen2025survey} and the ideal solution of size $12$ \cite{cmsv2024}, and genus pairs with $|U^*|+|V^*|=16,20,26,40$. So $f(A)\ge L+1$ already for $A/2\pi\ge2L-3$, $2\le L\le5$, and for $A/2\pi\ge10$, $18$ when $L=6,7$.
\end{example}

\begin{remark}[The first open case]\label{rem:T3}
Let $T_L$ be the least $|U^*|+|V^*|$ over genus-changing pairs sharing $L$ invariants; $T_2=6$ and $T_3\le10$. By Theorem~\ref{thm:descartes} a pair of size $8$ has shape $(3,5)$, and an exhaustive exact search excludes every one whose five-element side, scaled to coprime integers, has entries at most $220$ (supplement, Section~\sref{sec:searches}). Real solutions of that shape form a nonempty open set; for example $\{1,1,1,1,7\}$ has a partner triple of positive reals. So $T_3\in\{8,10\}$ is undecided, and an obstruction, if any, is arithmetic.
\end{remark}

\section{What heat does not hear}\label{sec:locality}

That the heat expansion cannot see the moduli is classical (Section~\ref{sec:prior}); we record it and then make the mechanism quantitative.

\begin{proposition}[Signature locality]\label{prop:locality}
For every $j\ge1$, $c_j(\Orb)=\Phi_j(\sigma(\Orb))$, where $\Phi_1(g;m)=\frac12(2g-2+\sum_i(1-\frac1{m_i}))$ and $\Phi_j(g;m)=\alpha_{j-1}\Phi_1(g;m)+\sum_ib_{j-2}(m_i)$ for $j\ge2$. In particular two orbifolds of the same signature have the same heat expansion to all orders.
\end{proposition}

\begin{proof}
This is Proposition~\ref{prop:heatinput} with \eqref{eq:area}: in Theorem~\ref{thm:IEH}, $\mathrm I$ depends only on the area, $\mathrm E$ only on the cone orders, and $\Hyp(t)=O(t^N)$ for every $N$ (Lemma~\ref{lem:hypbound}).
\end{proof}

\begin{proposition}[Moduli and rigidity]\label{prop:teich}\label{prop:rigidity}
For a closed orientable hyperbolic $2$-orbifold of signature $(g;m_1,\dots,m_n)$, $\Teich(\Orb)\cong\R^{6g-6+2n}$ \cite[Cor.~13.3.7]{thurston1980}; the dimension vanishes exactly for the triangle orbifolds. Two hyperbolic triangle orbifolds with the same signature are isometric, so $\Kiso(\Orb;\Sig_{0,3})=\Kmult(\Orb;\Sig_{0,3})$.
\end{proposition}

\begin{proof}
Thurston's dimension is $-3\chi(X_\Orb)+2k+l$ with $k=n$ cone points and $l=0$ corner reflectors, which vanishes only for $(g,n)\in\{(1,0),(0,3)\}$. For rigidity, let $\varphi$ be the M\"obius transformation taking the cone points of $\Orb$ to those of $\Orb'$, with equal orders corresponding. The pull-back of the metric of $\Orb'$ is a conformal metric of curvature $-1$ with cone angles $2\pi/m_i$, and by Troyanov's Theorem~A \cite{troyanov1991} (curvature $-1$, angles $2\pi/m_i$; its hypothesis of negative Euler characteristic is $\sum_i(1-1/m_i)>2$) such a metric is unique in its conformal class, so $\varphi$ is an isometry.
\end{proof}

\begin{corollary}[$\Kiso=\infty$ off the triangle orbifolds]\label{cor:Kinf}\label{prop:uncountable}
If $6g-6+2n>0$, the orbifolds of signature $\sigma=(g;m_1,\dots,m_n)$ form uncountably many isometry classes, and $\Kiso(\Orb;\Cl)=\infty$ for every class $\Cl$ containing them all (for example $\Sig$, or $\Sig_{0,n}$ when $g=0$, $n\ge4$) and every $\Orb$ of signature $\sigma$.
\end{corollary}

\begin{proof}
Points of $\Teich(\Orb)\cong\R^{6g-6+2n}$ are classes of discrete faithful representations of $\Gamma$ in $PSL(2,\R)$, and two of them give isometric orbifolds only if they differ by an automorphism of $\Gamma$ and a conjugation in $PGL(2,\R)$; so each isometry class meets $\Teich(\Orb)$ in an orbit of the countable group $\operatorname{Out}(\Gamma)$, and a countable union of countable sets is not $\R^d$, $d>0$. Any $\Orb'$ of signature $\sigma$ not isometric to $\Orb$ has $H_k(\Orb')=H_k(\Orb)$ for every $k$ (Proposition~\ref{prop:locality}).
\end{proof}

The shape is visible only beyond all orders, by the classical mechanism of Selberg, Huber and McKean \cite{huber1959,mckean1972,hejhal1976}: the heat traces of two orbifolds of one signature differ by the difference of their hyperbolic terms. We make this explicit: a bound with explicit constants and range (Theorem~\ref{thm:quantlocality}) and its sharpness: the leading term at the first length where the length spectra differ, so the factor $t^{-1/2}$ cannot be removed (Theorem~\ref{thm:sharp}, Corollary~\ref{cor:prefactor}).

\begin{theorem}[Quantitative locality]\label{thm:quantlocality}
Let $\Orb_1,\Orb_2\in\Sig$ have the same signature and area $A$, let $\ell$ be the smaller of their systoles (taken over all closed geodesics, including those through cone points) and $\diam$ the larger of their diameters.
\begin{enumerate}[label=(\alph*)]
\item For every $t>0$, $\cZ_{\Orb_1}(t)-\cZ_{\Orb_2}(t)=\Hyp_1(t)-\Hyp_2(t)$, so $|\cZ_{\Orb_1}(t)-\cZ_{\Orb_2}(t)|\le\max_i\Hyp_i(t)$.
\item For $0<t\le\ell^2/(2(1+\ell))$, with $B$ and $C$ as in Lemma~\ref{lem:hypbound} and \eqref{eq:Cconst},
\[
  |\cZ_{\Orb_1}(t)-\cZ_{\Orb_2}(t)|\le B(\ell,\diam,t)\le C(A,\ell,\diam)\,t^{-1/2}e^{-\ell^2/4t}.
\]
\end{enumerate}
\end{theorem}

\begin{proof}
(a) is Theorem~\ref{thm:IEH}, since $\mathrm I$ and $\mathrm E$ agree and $\Hyp_i\ge0$. (b) Lemma~\ref{lem:hypbound} bounds $\Hyp_i(t)$ by $B(\ell_i,\diam_i,t)$ whenever $t\le\ell_i^2/(2(1+\ell_i))$, which holds since $\ell\le\ell_i$ and $x\mapsto x^2/(2(1+x))$ is increasing. On the segment $[\ell,\ell_i]$ the function $B(\cdot,\diam_i,t)$ is decreasing, because $t\le x^2/(2(1+x))$ there, and $B$ is increasing in the diameter; so $B(\ell_i,\diam_i,t)\le B(\ell,\diam,t)$.
\end{proof}

The explicit constant is qualitative. It carries the factor $e^{3\diam}$ of Lemma~\ref{lem:counting}, through which alone the diameter, not determined by the signature, enters, and the admissible range ends where $e^{-\ell^2/4t}=e^{-(1+\ell)/2}$, so (b) says something only for small $t$ (the experiment in the supplement shows that it is loose); only the exponent and the factor $t^{-1/2}$ are sharp (Corollary~\ref{cor:prefactor}). For $x>0$ let $w_\Orb(x)=\sum_{[\gamma]:\,\ell(\gamma)=x}\ell(\gamma_0)$, the weighted length spectrum; it has finite support on every bounded interval.

\begin{theorem}[Sharpness]\label{thm:sharp}
Let $\Orb_1,\Orb_2\in\Sig$ have the same signature and write $w_i=w_{\Orb_i}$.
\begin{enumerate}[label=(\roman*)]
\item If $w_1\equiv w_2$, then $\cZ_{\Orb_1}\equiv\cZ_{\Orb_2}$ and the orbifolds are isospectral.
\item Otherwise let $L_*=\min\{x:w_1(x)\ne w_2(x)\}$. As $t\downarrow0$,
\[
  \cZ_{\Orb_1}(t)-\cZ_{\Orb_2}(t)=\frac{w_1(L_*)-w_2(L_*)}{2\sinh(L_*/2)}\,\frac{e^{-t/4}e^{-L_*^2/4t}}{\sqrt{4\pi t}}\big(1+o(1)\big).
\]
\item If $\ell_1\ne\ell_2$, then $L_*=\ell=\min(\ell_1,\ell_2)$.
\end{enumerate}
\end{theorem}

\begin{proof}
Group $\Hyp_1-\Hyp_2$ by length, and let $L'>L_*$ be the next length in the union of the supports of $w_1,w_2$. The terms of length at least $L'$ are bounded by the corresponding parts of $\Hyp_1+\Hyp_2$, which are $O(t^{-1/2}e^{-L'^2/4t})$ by the argument of Lemma~\ref{lem:hypbound} with $\ell$ replaced by $L'$; divided by the $L_*$ term this tends to $0$. If $w_1\equiv w_2$, the traces agree, and a heat trace determines the spectrum by uniqueness of the Laplace transform. For (iii), $w_i(x)=0$ for $x<\ell_i$ and $w_i(\ell_i)>0$.
\end{proof}

\begin{corollary}[What is attained]\label{cor:prefactor}
If the weighted length spectra first differ at the shorter systole $\ell$, in particular if the systoles differ, then $\sqrt t\,e^{\ell^2/4t}|\cZ_{\Orb_1}(t)-\cZ_{\Orb_2}(t)|\to|w_1(\ell)-w_2(\ell)|/(2\sinh(\ell/2)\sqrt{4\pi})\ne0$: the exponent $\ell^2/4$ and the factor $t^{-1/2}$ of Theorem~\ref{thm:quantlocality} are attained (the constant $C$ is not claimed to be). If they first differ at $L_*>\ell$, the difference is of the smaller order $t^{-1/2}e^{-L_*^2/4t}$, and if they never differ, it vanishes identically.
\end{corollary}

The factor $t^{-1/2}$ is the one-dimensional heat kernel along the closed geodesic: heat must travel the length $\ell$ in one dimension. Pairs with $\ell_1\ne\ell_2$ exist in every signature in which the systole varies on Teichm\"uller space; Figure~\ref{fig:F5}(b) shows them in the family $(0;3,3,3,3)$.

\begin{remark}[The expansion is blind; the trace is not]\label{rem:blind}
The heat trace as a function of $t$ determines the spectrum and the length spectrum \cite[Thm~1.1]{drydenstrohmaier2009}. For surfaces, length and eigenvalue spectra determine each other \cite[S\"atze~7--8]{huber1959}, the length spectrum determines a generic surface \cite{wolpert1979}, and isospectral sets are finite \cite{mckean1972}, as for hyperbolic orbisurfaces of genus at least one \cite[Thm~5.1]{dryden2004}. ``Heat does not hear the shape'' concerns the asymptotic series, not the spectrum.
\end{remark}

\section{The rigid case: triangle orbifolds}\label{sec:rigid}

A triangle orbifold $\Orb(p,q,r)$ is the double of a hyperbolic triangle with angles $\pi/p,\pi/q,\pi/r$, determined by its signature (Proposition~\ref{prop:rigidity}), so $\Kiso=\Kmult$ on $\Sig_{0,3}$. By \eqref{eq:s1inv} the first two heat invariants are equivalent to $\keytwo(p,q,r)=(S_1,R)$ and the third adds $P_3$. A \emph{collision} is a pair of distinct hyperbolic triads with the same $\keytwo$. We count prefixes $c_1,\dots,c_k$, the order in which the coefficients appear as $t\downarrow0$; on triangle orbifolds $c_2$ alone already gives $\keytwo$, since $12c_2+2=S_1+R$ with $S_1\in\Z$ and $0<R<1$. Off this family two invariants fail earlier, for instance on $(0;5,5,5)$ and $(0;2,2,2,10)$, of cone-order sums $15$ and $16$.

\begin{theorem}[Three heat invariants determine a triangle orbifold]\label{thm:three}
$R$, $S_1$ and $P_3$ determine $\{p,q,r\}$; hence $\Kiso(\Orb;\Sig_{0,3})\le3$, with equality exactly when $\Orb$ belongs to a collision, for instance $\Orb(2,8,8)$ and $\Orb(3,3,12)$.
\end{theorem}

\begin{proof}
With $e_1=S_1$, $e_2=Re_3$ and $P_3=e_1^3-3e_1e_2+3e_3$, $e_3=(P_3-e_1^3)/(3-3S_1R)$, where $S_1R\ge9$ by the arithmetic--harmonic mean inequality; the orders are the roots of $z^3-e_1z^2+e_2z-e_3$. The pair shares $\keytwo=(18,\frac34)$ and has $P_3=1032\ne1782$.
\end{proof}

This is Theorem~\ref{thmA} for $n=3$; isospectral hyperbolic triangle orbifolds are isometric, as also follows from \cite[Thm~1.1]{drydenstrohmaier2009}.

\subsection{Strata and the threshold}

A collision has equal sums, so injectivity of $\keytwo$ can be decided sum by sum. The hyperbolic triads $(p,q,r)$, $p\le q\le r$, with $p+q+r=S$ and least order $p$ form the \emph{$(S,p)$-stratum}; on it $R_{S,p}(q)=\frac1p+\frac1q+\frac1{S-p-q}$, $p\le q\le(S-p)/2$. For $p\ge3$ it is nonempty for $S\ge3p$ except $(S,p)=(9,3)$; for $p=2$ exactly when $S\ge11$, and it never contains $(2,2,S-4)$.

\begin{lemma}[Chamber monotonicity]\label{lem:chamber}
$R_{S,p}'(q)=-q^{-2}+(S-p-q)^{-2}\le0$, with equality only at $q=r$; so $R$ is injective on each stratum, with maximum at the spread triad $q=p$ (attained unless $p=2$ or $(S,p)=(9,3)$, whose spread triads have $R\ge1$) and minimum at the balanced triad $q=\lfloor(S-p)/2\rfloor$.
\end{lemma}

Let $R^+_{S,p}=\frac2p+\frac1{S-2p}$ and $R^-_{S,p}$ be these extreme values; $R^+_{S,p}$ is non-increasing in $p$ for $p\le S/3$, and since $\frac2{x-1}+\frac2{x+1}=\frac4x+\frac4{x(x^2-1)}$,
\begin{equation}\label{eq:Rminus}
  R^-_{S,p}=\frac1p+\frac4{S-p}+\frac{4\,[\,S-p\ \text{odd}\,]}{(S-p)((S-p)^2-1)} .
\end{equation}
For $p\ge2$ put $\varphi_p(S)=\frac4{S-p}-\frac1{S-2p-2}$, $\tau_p=\frac{p-1}{p(p+1)}$, $x^*(p)=\frac{3p(p+1)}{p-1}=3p+6+\frac6{p-1}$ and $\gap_p(S)=R^-_{S,p}-R^+_{S,p+1}$, which equals $\varphi_p(S)-\tau_p$ when $S-p$ is even and exceeds it when $S-p$ is odd. Two strata \emph{overlap} at $S$ if the convex hulls of their reciprocal sums meet.

\begin{lemma}[Adjacent strata]\label{lem:adjacent}
(i) If $\gap_p(S)>0$ whenever the strata $p$ and $p+1$ are both nonempty at $S$, then $\keytwo$ is injective on triads of sum $S$. (ii) For $S\ge3p+3$, the strata $p$ and $p+1$ overlap at $S$ if and only if $\gap_p(S)\le0$.
\end{lemma}

\begin{proof}
(i) For $p<p''$ the values of stratum $p''$ are at most $R^+_{S,p''}\le R^+_{S,p+1}<R^-_{S,p}$. (ii) For $S\ge3p+3$ the stratum $p+1$ contains its hyperbolic spread triad, so $R^-_{S,p}\le\max(\text{stratum }p+1)$ is $\gap_p(S)\le0$; the other condition, $R^-_{S,p+1}\le\max(\text{stratum }p)$, is automatic for $p\ge3$, and for $p=2$, where $\gap_2(S)\le0$ forces $S\ge18$ by (a), (b) below, $(2,5,S-7)$ has $R>\frac7{10}>\frac13+\frac{60}{224}\ge R^-_{S,3}$.
\end{proof}

\begin{theorem}[The first overlap]\label{thm:Sstar}
Let $S^*(2)=18$, $S^*(3)=19$, $S^*(p)=3p+8$ for $4\le p\le8$, and $S^*(p)=3p+7$ for $p\ge9$. For $S\ge3p+3$ the strata $p$ and $p+1$ overlap at $S$ if and only if $S\ge S^*(p)$.
\end{theorem}

\begin{proof}
(a) For $S-p$ even, $\varphi_p(S)-\tau_p=-\frac{(p-1)(S-3p-2)(S-x^*(p))}{p(p+1)(S-p)(S-2p-2)}$, which for $S\ge3p+3$ is $\le0$ exactly when $S\ge x^*(p)$. (b) For $S-p$ odd, $\gap_p>\varphi_p-\tau_p$, so nothing overlaps below $x^*(p)$; the numerator of $\varphi_p'$ is $-(S-3p-4)(3S-5p-4)$ and $x^*(p)>3p+4$, so the odd-parity gap decreases on $S\ge x^*(p)$. (c) $\gap_p(3p+7)=-\frac{2(p^2-5p-30)}{p(p+1)(p+3)(p+4)(p+5)}$ is negative exactly for $p\ge9$. (d) For $p\ge9$, $x^*(p)<3p+7$, which with (a)--(c) gives $S^*=3p+7$. For $2\le p\le8$, $x^*(p)=18,18,20,\frac{45}2,\frac{126}5,28,\frac{216}7$, the least $S\ge x^*(p)$ with $S-p$ even is $18,19,20,23,26,29,32$, the odd-parity sums in between ($S=18,28,31$ for $p=3,7,8$) have $\gap_p=\frac1{840},\frac1{2310},\frac1{10296}>0$, and at $S^*(p)+1$ the odd-parity gap is $-\frac7{936},-\frac1{72},-\frac{19}{3960},-\frac1{180},-\frac{76}{15015},-\frac1{231},-\frac{17}{4680}$.
\end{proof}

\begin{theorem}[Interval separation]\label{thm:separation}
For $S\le17$ the strata of sum $S$ have pairwise disjoint reciprocal sums. At $S=18$ the strata $2$ and $3$ meet in exactly $R=\frac34$, attained by $(2,8,8)$ and $(3,3,12)$, and all other strata are separated.
\end{theorem}

\begin{proof}
A collision joins two strata (Lemma~\ref{lem:chamber}), so some adjacent pair overlaps (Lemma~\ref{lem:adjacent}) and $S\ge S^*(p)\ge x^*(p)\ge18$, since $x^*(p)-18=3(p-2)(p-3)/(p-1)$. At $S=18$ only $S^*(2)\le18$, and $\gap_2(18)=0$ with $S-p$ even, so the strata $2,3$ share exactly $R^-_{18,2}=R^+_{18,3}=\frac34$; the strata $p\ge4$ lie below $R^+_{18,4}=\frac35$.
\end{proof}

\begin{theorem}[The threshold]\label{thm:threshold}
If $p+q+r\le17$, the first two heat invariants determine $\Orb(p,q,r)$ among all hyperbolic triangle orbifolds: $\Kiso(\Orb(p,q,r);\Sig_{0,3})\le2$. The bound is sharp (Theorem~\ref{thm:minimal}).
\end{theorem}

\begin{proof}
The first two invariants determine $S_1$ by \eqref{eq:s1inv}, so a triad of sum at most $17$ can collide only within its own sum, which Theorem~\ref{thm:separation} excludes.
\end{proof}

\begin{theorem}[The minimal degeneracy]\label{thm:minimal}
$\Orb(2,8,8)$ and $\Orb(3,3,12)$ share their first two heat invariants. They form the only collision of sum $18$, and no collision has a smaller sum.
\end{theorem}

\begin{proof}
$\frac12+\frac18+\frac18=\frac13+\frac13+\frac1{12}$ and $2+8+8=3+3+12$; the rest is Theorem~\ref{thm:separation}.
\end{proof}

Both theorems can be checked directly on the $83$ hyperbolic triads with $10\le S_1\le18$ (listed in the supplement); the strata explain the number $17$: at $S_1=18$ the most balanced order-$2$ and the most concentrated order-$3$ orbifold first reach the same total order and area (Figure~\ref{fig:F7}).

\paperfigure{F7}{fig:F7}

\begin{proposition}[Tangency only at $p=2$ and $p=4$]\label{prop:tangency}
For $S\ge3p+3$, the balanced triad of stratum $p$ and the spread triad of stratum $p+1$ have equal $R$ if and only if $p\in\{2,4\}$ and $S=x^*(p)$, that is, $(2,8,8),(3,3,12)$ at $S=18$ and $(4,8,8),(5,5,10)$ at $S=20$.
\end{proposition}

\begin{proof}
For $S-p$ even, equality forces $S=x^*(p)\in\Z$, so $(p-1)\mid6$, and of $p\in\{2,3,4,7\}$, with $x^*(p)-p=16,15,16,21$, only $p=2,4$ have the right parity. For $S-p$ odd, the decreasing odd-parity gap is positive at the last odd-parity sum below its zero and negative at the next (proof of Theorem~\ref{thm:Sstar}), so it vanishes at no such sum.
\end{proof}

Overlap is necessary but not sufficient: beyond it a collision needs two reciprocal sums from different strata to coincide exactly, an arithmetic event (the supplement tabulates first overlaps against first collisions for $p\le14$), and collisions also join non-adjacent strata, first $(5,15,15)$ and $(7,7,21)$ at $S=35$.

\begin{proposition}[Collision-free sums; computational]\label{prop:collisionfree}
Among the sums $18\le S\le4800$, exactly the $38$ sums
\[
\begin{gathered}
  19,\,21\text{--}25,\,27\text{--}30,\,33,\,41,\,44,\,46\text{--}51,\,59,\,65,\,67,\,81,\,99,\\
  115,\,119,\,123,\,125,\,173,\,199,\,203,\,223,\,235,\,243,\,251,\,307,\,329,\,557
\end{gathered}
\]
carry no collision. We do not know whether only finitely many sums are collision-free.
\end{proposition}

\begin{proof}[Method]
Exact rational arithmetic. (a) At each listed sum the reciprocal sums of all hyperbolic triads are pairwise distinct ($S=557$ has $25\,575$ triads). (b) Each of the other $4745$ sums carries an exhibited collision: $3962$ by scaling one at a divisor $d\mid S$, and $783$ by explicit pairs, the first being $(2,8,8),(3,3,12)$ at $18$, $(4,8,8),(5,5,10)$ at $20$ and $(4,10,12),(5,6,15)$ at $26$, listed in the supplementary data and each checked directly.
\end{proof}

\subsection{The minimal pair is isolated}\label{sec:isolation}

Because $S_1$ and $R$ scale oppositely, $\Lambda=S_1R$ is invariant under $(p,q,r)\mapsto(kp,kq,kr)$, and two triads with the same sum have the same $R$ exactly when they have the same $\Lambda$. A collision is therefore a set of positive rational points on the plane cubic
\[
  C_\Lambda:\ (X+Y+Z)(XY+YZ+ZX)=\Lambda\,XYZ,
\]
the curve that Bremner, Guy and Nowakowski study for integer $\Lambda$ \cite[p.~117]{bgn1993}. At $\Lambda=18\cdot\frac34=\frac{27}2$, which they do not treat, the minimal pair, $(1:4:4)$ and $(4:1:1)\sim(1:1:4)$, is a reciprocal pair in the sense of their remark on p.~117, which holds for every rational $\Lambda$.

\begin{theorem}[Isolation of the minimal pair]\label{thm:isolation}
The curve $C_{27/2}$ is a nonsingular cubic, and with base point $O=(1:-1:0)$ its group of rational points is finite of order $12$, isomorphic to $\Z/2\times\Z/6$. Its positive points are exactly the permutations of $(1:4:4)$ and $(1:1:4)$. Consequently, for every integer $k\ge1$ the orbifolds $\Orb(2k,8k,8k)$ and $\Orb(3k,3k,12k)$ share their first two heat invariants with each other and with no other hyperbolic triangle orbifold.
\end{theorem}

\begin{proof}
\emph{The model.} Put $e_2=XY+YZ+ZX$ and let
\begin{align*}
  \varphi(X:Y:Z)&=\Big(-\frac{16e_2}{Z^2},\ \frac{8(X-Y)}Z\Big(-\frac{4e_2}{Z^2}-54\Big)\Big),\\
  \psi(x,y)&=\big(25x+y:\ 25x-y:\ 4(x-216)\big).
\end{align*}
A direct substitution shows that $\psi$ maps $E:\ y^2=x(x+9)(x+384)=x^3+393x^2+3456x$ into $C_{27/2}$ and that $\varphi\circ\psi$ is the identity of $E$. The discriminant of $E$ is $2^{18}3^85^6\ne0$. The closure of $\psi(E)$ is an irreducible curve in the cubic $C_{27/2}$ birational to $E$, so of geometric genus $1$; it is not a line, a conic or a singular cubic, hence it is all of $C_{27/2}$, which is therefore a nonsingular cubic, and $\varphi,\psi$ are mutually inverse isomorphisms over $\Q$. On the three points with $Z=0$, $\varphi$ extends by $\varphi(O)=\infty$ (the origin of $E$), $\varphi(1:0:0)=(216,5400)$ and $\varphi(0:1:0)=(216,-5400)$; for example $\varphi(1:4:4)=(-24,360)$. So $C_{27/2}(\Q)\cong E(\Q)$.

\emph{Rank.} We use descent via $2$-isogeny \cite[\S3.6, Method~1 and (3.6.2)]{cremona1997}. Write $E:y^2=x(x^2+cx+d)$ with $c=393$, $d=3456=2^73^3$; the isogenous curve is $E':y^2=x(x^2+c'x+d')$ with $c'=-786$, $d'=c^2-4d=3^25^6$. For squarefree $d_1\mid d$ let $n_1$ be the number of classes $d_1$ for which $v^2=d_1u^4+cu^2+d/d_1$ has a rational point, and $n_1'$ likewise for $c',d'$; then $n_1=2^{e_1}$, $n_1'=2^{e_1'}$ and $\operatorname{rank}E(\Q)=e_1+e_1'-2$. A rational point gives coprime integers $M,e$ and an integer $N$ with $N^2=d_1M^4+cM^2e^2+(d/d_1)e^4$.
(i) On $E$, $d_1\in\{\pm1,\pm2,\pm3,\pm6\}$. For $d_1\in\{\pm2,\pm3\}$, $d_1\equiv\pm2$ is a non-residue modulo $5$, $d/d_1\equiv d_1^{-1}$ (as $d\equiv1$), $d_1+d_1^{-1}\equiv0$ and $c\equiv3$. Since $M^4,e^4\in\{0,1\}$ modulo $5$ and $5$ divides at most one of $M,e$, the right side is $\equiv d_1$ if $5\mid e$, $\equiv d_1^{-1}$ if $5\mid M$, and $\equiv d_1+d_1^{-1}+3M^2e^2\equiv3M^2e^2\equiv\pm3$ if $5\nmid Me$, each a non-residue. The classes $\pm1,\pm6$ are the images of $O,(0,0),(-9,0),(-384,0)$. So $n_1=4$, $e_1=2$.
(ii) On $E'$, $d_1\in\{\pm1,\pm3,\pm5,\pm15\}$. If $d_1<0$ all three coefficients are negative. For $d_1=3,5,15$ write $f=d_1M^4-786M^2e^2+(d'/d_1)e^4=3^vf_0$ with $3\nmid f_0$; then $v$ is odd or $f_0\equiv2\pmod3$: for $d_1=3$, $f=3(M^4-262M^2e^2+15625e^4)$ with bracket $\equiv1\pmod3$; for $d_1=5$, $f\equiv2M^4\pmod3$, and if $M=3M_1$ then $f=9(45M_1^4-786M_1^2e^2+3125e^4)$ with bracket $\equiv2$; for $d_1=15$, $f=3g$ with $g\equiv2M^4-M^2e^2+2e^4\pmod3$, so $g\equiv2$ if $3$ divides $M$ or $e$, and otherwise $g\equiv e^4(5w^2-w+2)\equiv6e^4\pmod9$ with $w\equiv M^2/e^2\in\{1,4,7\}$, so $f=9g'$ with $g'\equiv2\pmod3$. So $n_1'=1$, $e_1'=0$, and the rank is $0$; the Selmer groups are $\{\pm1,\pm6\}$ and $\{1\}$, every class being decided by a local condition. Independently, PARI/GP's \texttt{ellrank} returns rank $0$ for $E$ (Appendix~\ref{app:reproducibility}).

\emph{Torsion.} $E$ has good reduction at $7$ and $11$, the reduction map is injective on torsion at an odd prime of good reduction \cite[\S3.3, p.~70]{cremona1997}, and $\#E(\mathbb F_7)=\#E(\mathbb F_{11})=12$. The twelve points $\infty$, $(0,0)$, $(-9,0)$, $(-384,0)$, $(16,\pm400)$, $(-24,\pm360)$, $(-144,\pm2160)$, $(216,\pm5400)$ lie on $E$: the three points with $y=0$ have order $2$; $(16,400)$ is a flex, since the tangent $y=21x+64$ meets $E$ only there, with multiplicity $3$, so $2(16,400)=(16,-400)=-(16,400)$ and $(16,\pm400)$ have order $3$; and the remaining six have order $6$. So $E(\Q)\cong\Z/2\times\Z/6$.

\emph{The points of $C_{27/2}$.} The twelve points $(1:0:0)$, $(0:1:0)$, $(0:0:1)$, $(1:-1:0)$, $(0:1:-1)$, $(1:0:-1)$ and the permutations of $(1:4:4)$ and $(1:1:4)$ lie on $C_{27/2}$ and are distinct, so they are all of $C_{27/2}(\Q)$, and the positive ones are the six permutations. A triad with $S_1=18k$ and $R=\frac3{4k}$ has $\Lambda=\frac{27}2$, so it is a permutation of $(1:4:4)$ or $(1:1:4)$, and the only integer representatives with sum $18k$ are $(2k,8k,8k)$ and $(3k,3k,12k)$; both are hyperbolic, since $R<1$ and all entries are at least $2k\ge2$.
\end{proof}

\section{Stability}\label{sec:stability}

We study the conditioning of the algebraic inverse problem from $(c_1,\dots,c_n)$ to the positive real orders $m$ of a sphere with $n$ cone points, $\mu=\max_im_i$; for real $m$ the $c_j$ are the rational functions of Proposition~\ref{prop:heatinput} evaluated at $m$, heat invariants of an orbifold only for integer orders. We proceed through its three layers: a triangular linear map, the system of Theorem~\ref{thmB}, and polynomial roots.

\begin{remark}[What is measured]\label{rem:stabscope}
The error model is on $c_1,\dots,c_n$, not on eigenvalues, and genus $0$ and a known $n$ are assumed. The constants are evaluated at the true orders, unknown in an application, so the logic is a-posteriori: recover a candidate, then certify it with Proposition~\sref{prop:S5} of the supplement at the candidate, with radii covering the distance from the data to its exact invariants. The certificate is as reliable as those radii; in the blind recovery of the supplement they come from a-posteriori error bars, so there it validates the recovery but does not certify it.
\end{remark}

The \emph{recovery map} is: (1) $\tilde\cI=\mathbf F^{-1}(\tilde c-h_0)$ (Proposition~\ref{prop:frontend}); (2) $\tilde e=M(\tilde\cI)^{-1}b(\tilde\cI)$ (Theorem~\ref{thmB}); (3) the roots of $\tilde q(z)=\sum_j(-1)^j\tilde e_jz^{n-j}$, $\tilde e_0=1$; (4) for integer orders, rounding their real parts. Hats denote scale-free quantities, $\hat m=m/\mu$, $\hat e_k=e_k/\mu^k$, $\hat\cI=(\mu R,P_1/\mu,P_3/\mu^3,\dots)$, and $\hat e$ solves $M(\hat\cI)\hat e=b(\hat\cI)$ because the system is weighted homogeneous. Distances between multisets are $\dist(m,m')=\min_\pi\max_i|m_i-m'_{\pi(i)}|$.

\begin{proposition}[Front end]\label{prop:frontend}
With $p_\nu(x)=\sum_ka_{\nu,k}x^{2k}$, $(c_1,\dots,c_n)^\top=\mathbf F\,\cI_n+h_0(n)$, where $\mathbf F$ is lower triangular, indexed from $0$, with $\mathbf F_{00}=-\frac12$, $\mathbf F_{\nu+1,0}=-\frac12\alpha_{\nu+1}+(-1)^\nu a_{\nu,0}$, $\mathbf F_{\nu+1,k}=(-1)^\nu a_{\nu,k}$ ($1\le k\le\nu+1$), and $h_0(n)=\frac{n-2}2(1,\alpha_1,\dots,\alpha_{n-1})$. $\mathbf F$ does not depend on $n$, and its diagonal entries $(-1)^\nu|B_{2\nu+2}|/(2(\nu+1)!(2\nu+1))$ are nonzero.
\end{proposition}

\begin{proof}
By \eqref{eq:cj}, $c_1=(n-2-R)/2$ and $c_{\nu+2}=\alpha_{\nu+1}(n-2-R)/2+(-1)^\nu\sum_i\sum_ka_{\nu,k}m_i^{2k-1}$; the diagonal is Lemma~\ref{lem:conepoly}.
\end{proof}

If every $|\delta c_j|\le\delta$, the $r$-th invariant has error at most $\amp_r\delta$, the absolute row sum of $\mathbf F^{-1}$: $\amp_0,\dots,\amp_4=2,14,498,4062,\frac{56230}3$ (for example $P_3=-18\tilde c_1-120\tilde c_2-360\tilde c_3$, $\tilde c=c-h_0$).

\begin{lemma}[Hurwitz factorisation]\label{lem:hurwitz}
Let $\mathbf B,\mathbf Y$ be the $n\times n$ matrices with rows $0\le k\le n-1$, $\mathbf B_{k,j}=(-1)^{j+1}e_{2k+1-j}$, $\mathbf Y_{k,i}=e_{2(k-i)}$ for $i\le k\le n-2$, $\mathbf Y_{n-1,n-1}=(-1)^ne_n$, other entries of $\mathbf Y$ zero ($e_i=0$ outside $[0,n]$). Then $\mathbf B=\mathbf YM$, $\det\mathbf B=(-1)^{n(n-1)/2}\prod_{i<j}(m_i+m_j)$ and $M^{-1}=\mathbf B^{-1}\mathbf Y$.
\end{lemma}

\begin{proof}
For $f=\prod_i(1+m_iz)=E_f(z^2)+zO_f(z^2)$ and $d(z)=\sum_kd_kz^k=E_d(z^2)+zO_d(z^2)$, $\mathbf Bd$ is the coefficient vector of $E_fO_d-E_dO_f$, and $(Md)_j=[z^{2j+1}](zO_d-TE_d)$ for $j\le n-2$; since $T\equiv zO_f/E_f$ modulo $z^{2n-1}$, comparing coefficients of $E_f(zO_d-TE_d)$ gives the rows $k\le n-2$, and the top row uses $R=e_{n-1}/e_n$. Then $\det\mathbf Y=(-1)^ne_n$ and Theorem~\ref{thmB}.
\end{proof}

So inverting Theorem~\ref{thmB} is as well conditioned as the Hurwitz matrix of $\chi_m$. Let $\cond=\|M(\hat\cI)^{-1}\|_\infty$, $\mathfrak t_k(n)=[w^k]\operatorname{artanh}(w)\sec^2((n+1)\operatorname{artanh}w)$, $\zeta_n=\max\big(1,\max_j\sum_{0\le i<j,\,2\le2j-2i\le n}\mathfrak t_{2i+1}(n)\big)$ (so $\zeta_3=1$, $\zeta_4=\frac{79}3$, $\zeta_5=\frac{14048}{15}$) and $\mathfrak r_n(\hat m)=\max\big(\hat e_n,\max_{j\le n-2}\sum_{i\le j}\mathfrak t_{2i+1}(n)\hat e_{2j-2i}\big)$.

\begin{theorem}[Lipschitz dependence of $e$ on $\cI_n$]\label{thm:S2}
Let $\eta=\max\big(\mu|\tilde R-R|,\max_{1\le l\le n-1}|\tilde P_{2l-1}-P_{2l-1}|/\mu^{2l-1}\big)$ and assume $\eta\le1$. (a) $\cond\le n\binom{2n}n^{n/2}2^{n-1}/\prod_{i<j}(\hat m_i+\hat m_j)$. (b) If $\cond\,\zeta_n\eta\le\frac12$, then $M(\tilde\cI)$ is invertible and $\max_k|\tilde e_k-e_k|/\mu^k\le2\cond\,\mathfrak r_n(\hat m)\,\eta$. (c) To first order $\tilde e-e=-M^{-1}J(\tilde\cI-\cI)$, with $J$ as in the proof of Theorem~\ref{thmB}.
\end{theorem}

\begin{proof}
(a) Lemma~\ref{lem:hurwitz}, Hadamard's inequality for the cofactors of $\hat{\mathbf B}$ (whose columns contain $\hat e_0=1$ or $\hat e_1\ge1$), $|\det\hat{\mathbf B}|=\prod(\hat m_i+\hat m_j)$ and $\hat e_k\le\binom nk$. (b) All $P_k>0$, so coefficientwise $\hat U\le n\operatorname{artanh}w$ and $|\delta\hat U|\le\eta\operatorname{artanh}w$; as the coefficients of $\tanh$ are those of $\tan$ up to sign, $|\tilde T_k-T_k|\le\eta\,\mathfrak t_k(n)$. Hence the residual $r=b(\tilde\cI)-M(\tilde\cI)e$ has $\|r\|_\infty\le\mathfrak r_n\eta$, $\|M(\tilde\cI)-M(\cI)\|_\infty\le\zeta_n\eta$, and $M(\tilde\cI)(\tilde e-e)=r$. (c) is the derivative computed for Theorem~\ref{thmB}.
\end{proof}

On every test multiset the exact $\cond$ is at most $3.511$; the Hadamard step behind (a) is the lossy one. Ostrowski's global root bound \cite[Th\'eor\`eme~XXX, (71,1)]{ostrowski1940} has exponent $1/n$ everywhere; Rouch\'e's theorem gives the local one.

\begin{theorem}[Heat invariants to cone orders]\label{thm:S3}
Let $\delta=\max_j|\tilde c_j-c_j(m)|$, $\Xi_\mu=\max(\mu\amp_0,\max_{r\ge1}\amp_r\mu^{1-2r})$, and suppose $\Xi_\mu\delta\le1$ and $\cond\,\zeta_n\Xi_\mu\delta\le\frac12$. For a distinct order $a$ of multiplicity $k_a$, with $\widehat{\sepn}_a=\min_{b\ne a}|\hat a-\hat b|$ and $\hat\Pi_a=\prod_{b\ne a}|\hat a-\hat b|^{k_b}$, if
\[
  r_a=\mu\Big(\frac{2^{2-k_a}3^n\cond\,\mathfrak r_n\,\Xi_\mu\,\delta}{\hat\Pi_a}\Big)^{1/k_a}\le\mu\,\frac{\min(\widehat{\sepn}_a,1)}2,
\]
then the recovered polynomial, rescaled, has exactly $k_a$ roots within $r_a$ of $a$. So $\dist(\text{roots},m)\le\max_aC_a\delta^{1/k_a}$ with $C_a=r_a\delta^{-1/k_a}$ when this holds for every $a$.
\end{theorem}

\begin{proof}
Proposition~\ref{prop:frontend} gives $\eta\le\Xi_\mu\delta$, and Theorem~\ref{thm:S2}(b) gives $|\tilde e_k-\hat e_k|\le\varepsilon=2\cond\,\mathfrak r_n\Xi_\mu\delta$. On $|z-\hat a|=r\le\min(\widehat{\sepn}_a,1)/2$, $|q(z)|\ge r^{k_a}\hat\Pi_a2^{-(n-k_a)}$ while $|z|\le\frac32$ gives $|\tilde q-q|<2^{1-n}3^n\varepsilon$; apply Rouch\'e's theorem.
\end{proof}

So recovery is Lipschitz at simple orders and H\"older of exponent $1/k$ at a $k$-fold coincidence. The exponent $1/k$ is sharp for arbitrary data (Proposition~\ref{prop:sharpexp}(ii)); for data coming from real orders it is $\frac12$ at a double order (Proposition~\ref{prop:sharpexp}(i)) and when all the orders are equal (Remark~\ref{rem:triple}); other configurations are not settled (Figure~\ref{fig:F8}).

\begin{proposition}[Sharpness of the exponents]\label{prop:sharpexp}
(i) \emph{Double order, realisable data.} For $m(s)=(a+s,a-s,m_3,\dots,m_n)$ the heat invariants are smooth and even in $s$, so they move by $O(s^2)$ while the orders move by $s$; for $(2,8,8)$, $\delta R=s^2/(4(64-s^2))$, $\delta P_1=0$, $\delta P_3=48s^2$. (ii) \emph{$k$-fold order, general data.} Let $q_s(z)=((z-a)^k-s^k)g(z)$, $g$ real and monic of degree $n-k$, $a\ne0$, $g(0)\ne0$, $\prod_{i<j}(z_i+z_j)\ne0$ for the roots of $q_0$. For small $s>0$ the data of $q_s$ are within $O(s^k)$ of those of $q_0$, the system of Theorem~\ref{thmB} is invertible at them and returns $q_s$, and its roots $a+se^{2\pi ij/k}$ lie at distance $s$ from $a$. For $k\ge3$ these roots are not all real, so these data are not the heat invariants of any real multiset.
\end{proposition}

\begin{proof}
(i) Exchanging $a\pm s$ fixes the multiset. (ii) The data are smooth in the coefficients near $q_0$, as $e_n(q_0)\ne0$, and invertibility persists for small $s$; a real multiset with these data would have the elementary symmetric functions of $q_s$.
\end{proof}

\begin{remark}[Realisable data when all orders are equal]\label{rem:triple}
For data that are heat invariants of real orders near $(a,\dots,a)$, $n\ge3$, the exponent is $\frac12$: for real $d_i\ge-a$ (in particular $|d_i|\le a$), $(P_3-3a^2P_1)(a+d)-(P_3-3a^2P_1)(a)=\sum_id_i^2(3a+d_i)\ge2a\sum_id_i^2$, so under $|\delta c_j|\le\delta$, $\max_i|d_i|\le((498+42a^2)\delta/(2a))^{1/2}$. The family $(a+s,a-s,a,\dots,a)$, with $\delta P_1=0$, $\delta P_3=6as^2$ and all other changes $O(s^2)$, attains $\frac12$. For $n=1$ recovery is Lipschitz; mixed clusters are not settled.
\end{remark}

\paperfigure{F8}{fig:F8}

\begin{theorem}[Explicit threshold]\label{thm:S4}
For integer orders $m$ let
\[
  \delta_{\rm thm}(m)=\min\Big(\frac1{\Xi_\mu},\ \frac1{2\cond\,\zeta_n\Xi_\mu},\ \min_a\frac{\hat\Pi_a\,2^{k_a-1}}{3^n(2\mu)^{k_a}\cdot2\cond\,\mathfrak r_n\Xi_\mu}\Big).
\]
If $|\tilde c_j-c_j(m)|\le\delta_{\rm thm}(m)$ for all $j$, rounding the real parts of the recovered roots returns $m$.
\end{theorem}

\begin{proof}
The third term makes $r_a\le\frac12$ in Theorem~\ref{thm:S3}; the discs $|z-a|<\frac12$ are disjoint and contain $k_a$ roots each, hence all of them.
\end{proof}

A sharper certificate, equally rigorous for errors in the heat invariants, runs the same three layers with coefficientwise bounds in exact rational arithmetic (Proposition~\sref{prop:S5} of the supplement); it gives $\delta_{\rm cert}$ in Table~\ref{tab:thresholds}. The failure bound $\delta_{\rm up}$ minimizes $\|c(\tilde q)-c(m)\|_\infty$ over real monic $\tilde q$ with a root, or a complex pair, of real part $a\pm\frac12$, where rounding is a tie.

% BEGIN GENERATED TABLE thresholds
\begin{table}[t]
\caption{Thresholds for exact recovery of integer orders under the uniform model $|\delta c_j|\le\delta$ on the heat invariants: the closed form $\delta_{\rm thm}$ of Theorem~\ref{thm:S4}, the threshold $\delta_{\rm cert}$ certified in exact arithmetic by Proposition~\sref{prop:S5} of the supplement, and a constructed failure $\delta_{\rm up}$; aeb means $a\times10^b$, $\delta_{\rm thm}$ and $\delta_{\rm cert}$ are rounded down and $\delta_{\rm up}$ up. Table~\sref{tab:thresholdsfull} adds the relative precisions.}\label{tab:thresholds}
\centering\footnotesize
\begin{tabular}{@{}llll@{}}
\toprule
$m$ & $\delta_{\rm thm}$ & $\delta_{\rm cert}$ & $\delta_{\rm up}$\\
\midrule
$(2,8,8)$ & $3.80\text{e}{-7}$ & $2.341\text{e}{-3}$ & $2.485\text{e}{-3}$ \\
$(3,3,12)$ & $1.18\text{e}{-7}$ & $4.040\text{e}{-3}$ & $4.589\text{e}{-3}$ \\
$(3,10,15,30)$ & $4.02\text{e}{-11}$ & $3.660\text{e}{-3}$ & $7.487\text{e}{-3}$ \\
$(4,5,21,28)$ & $3.14\text{e}{-11}$ & $1.461\text{e}{-3}$ & $2.018\text{e}{-3}$ \\
$(2,3,7)$ & $4.49\text{e}{-7}$ & $3.658\text{e}{-3}$ & $6.587\text{e}{-3}$ \\
$(4,4,4)$ & $9.35\text{e}{-7}$ & $4.539\text{e}{-4}$ & $5.036\text{e}{-4}$ \\
$(7,7,7)$ & $9.97\text{e}{-8}$ & $8.068\text{e}{-5}$ & $8.273\text{e}{-5}$ \\
$(3,3,4,4)$ & $1.48\text{e}{-9}$ & $9.597\text{e}{-5}$ & $1.195\text{e}{-4}$ \\
$(5,5,5,5)$ & $4.74\text{e}{-10}$ & $3.617\text{e}{-5}$ & $3.826\text{e}{-5}$ \\
$(2,2,2,3)$ & $2.03\text{e}{-9}$ & $1.858\text{e}{-4}$ & $2.520\text{e}{-4}$ \\
$(2,2,2,2,3)$ & $2.72\text{e}{-12}$ & $7.908\text{e}{-6}$ & $5.312\text{e}{-5}$ \\
\bottomrule
\end{tabular}
\end{table}
% END GENERATED TABLE thresholds

The certified threshold is within a factor $1.03$--$2.05$ of a constructed failure for $n\le4$ and $6.72$ for $n=5$: close to the truth, not just safe (only the exponents are claimed sharp).

\section{Numerical illustration}\label{sec:numerics}

Nothing here is used in a proof; the full experiments, a blind recovery and the moduli figure are in the supplement. The spectrum of $\Orb(p,q,r)$ is the union of the Neumann and Dirichlet spectra of its triangle (the identity U\c{c}ar uses for lunes \cite[Thm~4.10, Cor.~4.18]{ucar2017}), computed with $H^1$ finite elements of order $10$ on meshes of hyperbolic size $0.05$ (NETGEN \cite{schoberl1997}, NGSolve \cite{ngsolve}) and shift-invert Lanczos iteration \cite{arpack1998}: about $2850$ eigenvalues per orbifold, up to $\lambda\approx2.17\times10^4$ (Neumann) and $2.38\times10^4$ (Dirichlet). Reflection makes the eigenfunctions smooth at the corners, and convergence is fast ($h^{13}$ to $h^{19}$ at order $10$).

\emph{Error budget.} The error estimate of each eigenvalue is the difference between two independent discretisations, for the first $500$ at most $2.9\times10^{-11}$ relative. With the mean value theorem and a tail bound from the Weyl law it bounds the error of the difference of the two traces on $0.0015\le t\le0.012$ by $2.9\times10^{-11}$: an a-posteriori estimate, not a certified enclosure. Each trace agrees with $\mathrm I+\mathrm E$ to $7\times10^{-13}$ for $0.0015\le t\le0.03$, so no eigenvalue below about $1.6\times10^4$ is missing; above that the test is blind, and a double-window recomputation of the spectra has not been run. The method, a Bolza benchmark and these limitations are detailed in Section~\sref{sec:listening} of the supplement.

\emph{Two timescales.} $D(t)=\cZ_{\Orb(2,8,8)}(t)-\cZ_{\Orb(3,3,12)}(t)=\sum_{j\ge3}d_jt^{j-2}$ (a function of $t$, not the diameter) has $d_3=\frac{25}{12}$, $d_4=-\frac{1775}{24}$, $d_5=\frac{153025}{48}$ by Proposition~\ref{prop:heatinput}. An unweighted least-squares fit of $D(t)/t$ by a polynomial of degree $9$ on $[0.0015,0.012]$ (window and degree chosen blind, minimizing the larger of the data error and the change on adding a term) gives $d_3=2.0833320\pm0.0000031$ and $d_4=-73.9551\pm0.0062$, within $0.42$ and $0.53$ of these heuristic error bars of the exact values. Pointwise, for $0.0015\le t\le0.03$, $D(t)$ agrees with the exact difference of the elliptic terms to $4\times10^{-13}$, testing every order of the cone expansion at once. The window is small because the cone series diverges, its coefficients growing like $l!\,(144/\pi^2)^l$ (the poles of $\Phi_{12}$ at $\pm\pi/12$): at $t=0.02$, $d_3t,d_4t^2,d_5t^3$ are $0.042,-0.030,0.026$. At larger $t$ the closed geodesics take over; fitting $\cZ-\mathrm I-\mathrm E$ on $[0.06,0.5]$ gives leading multiplicities $1.997$ and $1.9998$ at the shortest lengths $2.2568$ and $1.8626$: one geodesic, two orientations (Figure~\ref{fig:F5}(a)).

\paperfigure{F5}{fig:F5}

\emph{Changing the shape.} Eight non-isometric members of a family of signature $(0;3,3,3,3)$ have heat traces equal to $\mathrm I+4\mathrm E_3$ to $1.2\times10^{-12}$ where geodesics are negligible, while $\lambda_1$ falls from $4.122$ to $0.435$; for all $28$ pairs the difference of traces is the difference of hyperbolic terms, computed with nothing fitted, and its ratio to the leading term of Theorem~\ref{thm:sharp} approaches $1$ (Figure~\ref{fig:F5}(b); supplement, Section~\sref{sec:moduli}).

\section{Open problems}\label{sec:open}

\begin{problem}[Linear growth]\label{prob:growth}
Is $N(k)=O(k)$, that is (Theorem~\ref{thm:growth}), is $f(A)=\Theta(A)$?
\end{problem}

\begin{problem}[The first genus collision]\label{prob:T3}
Is $T_3=8$ or $10$ (Remark~\ref{rem:T3})?\end{problem}

\begin{problem}[Integer sharpness]\label{prob:sharp}
For $n\ge5$, do hyperbolic spheres with $n$ cone points and different orders share $n-1$ heat invariants? By Theorem~\ref{thmC}(1) this asks for an odd symmetric Prouhet--Tarry--Escott system \cite[\S3]{borweiningalls1994} of $n$ positive and $n$ negative integers with vanishing reciprocal sum; yes for $n=3,4$, and over the reals.
\end{problem}

\backmatter

\bmhead{Acknowledgements}
The colour schemes of the figures are Scientific Colour Maps \cite{crameri2023}, designed to be perceptually uniform and readable by people with colour-vision deficiencies \cite{crameri2020}.

\section*{Statements and Declarations}

\bmhead{Funding}
No funding was received for this work.

\bmhead{Competing interests}
The authors have no competing interests to declare that are relevant to the content of this article.

\bmhead{Ethics approval and consent}
Not applicable.

\bmhead{Data and code availability}
All code and data are in the public repository \url{https://github.com/Ali-M658/Arithmetic-verification}. One command, \texttt{bash code/run\_all.sh}, reruns every check (with \texttt{--quick} for the exact-arithmetic suite and \texttt{--full} for the long computations), and \texttt{DATA-MANIFEST.md} lists every data file with its generating script and checksum (Appendix~\ref{app:reproducibility}). An archived copy will be deposited at Zenodo: \textbf{[PLACEHOLDER: Zenodo DOI, to be minted at submission.]}

\bmhead{Author contributions}
\textbf{[PLACEHOLDER: author contribution statement, to be supplied by the authors before submission.]}

\begin{appendices}

\section{Proofs of the growth results}\label{app:pte}

\begin{lemma}[Pigeonhole]\label{lem:pigeon}
For $L\ge2$ there are distinct $n$-multisets of positive integers with equal power sums of every odd exponent $j\le2L-3$, with $n\le(L-1)^2+1$, and with $n\le N(2L-3)$.
\end{lemma}

\begin{proof}
For $n=(L-1)^2+1$ and $M>n!\,n^{L-1}$, the $\ge M^n/n!$ multisets of size $n$ in $\{1,\dots,M\}$ give at most $n^{L-1}M^{(L-1)^2}<M^n/n!$ vectors $(P_1,P_3,\dots,P_{2L-3})$; for the second bound translate a solution of degree $2L-3$.
\end{proof}

\begin{proposition}[Shift]\label{prop:shift}
Let $X\ne Y$ be disjoint $n$-multisets of integers, $n\ge1$, with equal power sums of every exponent $j\le k$, $k\ge2L-3$, and for $c\in\Q\setminus(-X\cup-Y)$ let $Z(c)=(X+c)\uplus(-(Y+c))$. Then (1) $\sum_{Z(c)}z^j=0$ for odd $j\le2L-3$; (2) $\iota(Z(c))=2(\#\{x>-c\}-\#\{y>-c\})$ is $0$ for $c>-\min(X\cup Y)$ and nonzero on an open interval; (3) $\varrho(c)=\sum_{Z(c)}z^{-1}$ is a nonzero rational function. None of these changes when pairs $\{z,-z\}$ are removed.
\end{proposition}

\begin{proof}
(1) $P_j(X+c)-P_j(Y+c)$ is a combination of the $P_i(X)-P_i(Y)$, $i\le j$. (2) Count signs; the counting functions of $X\ne Y$ differ on an interval. (3) The poles are distinct and $n\ge1$.
\end{proof}

\begin{proof}[Proof of Theorem~\ref{thm:growth}]
(a) Lemma~\ref{lem:pigeon} and Proposition~\ref{prop:doubling} give two genus-$0$ orbifolds of different signatures and area below $2\pi(3(L-1)^2+1)$ sharing $H_L$, so with $\Kmult\ge L+1$; take the largest $L\ge2$ with $3(L-1)^2+1\le A/2\pi$. The upper bound is Corollary~\ref{cor:sigarea}.
(b) A partner sharing $H_L$ gives $Z$ with $|Z|\le|U|+|V|\le2\lfloor A/\pi\rfloor+8$ (Corollary~\ref{cor:sigarea}), and $Z,-Z$ solve degree $2L-2$.
(c) By Lemma~\ref{lem:pigeon} the pair in (a) can have $n\le N(2L-3)\le C(2L)^\beta$, so $f(A)\ge L+1$ once $6\pi C(2L)^\beta\le A$; take $L=\lfloor\frac12(A/6\pi C)^{1/\beta}\rfloor$ if $L\ge2$, and otherwise use $f(A)\ge3$ ($\Orb(2,8,8)$ has area $\pi/2$ and $\Kmult=3$).
(d) ``If'' is (c). Conversely, $N$ is nondecreasing, and $L=\lceil k/2\rceil+1$, $A=((L+1)/c)^{1/\alpha}$ in (b) give $N(k)\le2A/\pi+8=O(k^{1/\alpha})$.
(e) \emph{Genus.} Take $X,Y$ attaining $N(2L-3)$, common elements cancelled; by Proposition~\ref{prop:shift} choose rational $c$ with $\iota(Z(c))\ne0\ne\varrho(c)$ and $c'$ large with $\iota(Z(c'))=0\ne\varrho(c')$. Then $Z=Z(c)\uplus\lambda Z(c')$, $\lambda=-\varrho(c')/\varrho(c)$, is after removing pairs an $L$-configuration with $\iota(Z)\ne0$ and $|Z|\le4N(2L-3)$, realised (Section~\ref{sec:pte}) with different genera; the side $W$ with more entries has the smaller genus $g_0$, and the area is below $2\pi(|W|-2)$ if $g_0=0$ and at most $4\pi$ if $g_0=1$, so below $2\pi|Z|\le8\pi N(2L-3)$. \emph{Cone count.} With $X,Y$ translated so that $1\in X$ is least, Proposition~\ref{prop:doubling} gives cone counts $3n-\mult_X\ne3n$ and area below $2\pi(3N(2L-3)-2)$. As $N(2L-3)\le2L^2$, $f_g(A)\ge L+1$ for $A\ge16\pi L^2$ and $f_n(A)\ge L+1$ for $A\ge12\pi L^2$, giving the bounds for $A\ge64\pi$, $48\pi$; below, $f_g\ge3$ by Example~\ref{ex:siggenus}, and $f_n\ge3$ as $(0;5,5,5)$, $(0;2,2,2,10)$ (area $2\pi\cdot\frac25$) share two invariants. (d) for $f_g,f_n$ follows from these thresholds and $f_g,f_n\le f$.
\end{proof}

\section{Computational methods and reproducibility}\label{app:reproducibility}

\bmhead{Use of AI tools}
\textbf{[PLACEHOLDER: statement on the use of AI tools in preparing this manuscript, to be written by the authors. The journal's guidelines ask for any such use to be documented in the Methods section or, where there is none, in a suitable alternative part of the manuscript; this appendix serves as that part.]}

\bmhead{Computations}
Heat invariants, certificates and search results are exact (integer or rational arithmetic) unless stated, and each is produced by a committed program of the repository named in the data statement (paths from its root). (i)~Remark~\ref{rem:witnesses}, $n=4$: all pairs of $4$-multisets of orders in $[2,440]$ with equal $(R,P_1,P_3)$, found by sorting every multiset on these exact keys ($193$ pairs, $107$ primitive; \repo{review/audit-2/pte-witnesses/check_direct4.c}); pencil splittings: \repo{review/round1-fixes/d1_pencil_counts.py}. $n=5$: all $216{,}071{,}394$ multisets of orders in $[2,120]$, grouped by $P_1$ and keyed on the exact integers $(P_3,P_5,\operatorname{lcm}(2,\dots,120)\,R)$ (\repo{theory/audibility/sharpness_search.py}). (ii)~Remark~\ref{rem:T3}: for each $5$-multiset $V$ with gcd $1$ and entries at most $220$, the other side is the root set of an explicit cubic; a C filter rejects $V$ only when floating-point inclusion disks with rigorous error bounds exclude three rational roots, and every other $V$ is decided exactly (\repo{theory/pte/search_T3.c}, \repo{search_T3_verify.py}, log \repo{theory/pte/data/T3_search_log.txt}, planted controls in \repo{search_T3_control.py}). (iii)~Proposition~\ref{prop:collisionfree}: the triads of each listed sum, and an exhibited pair at each other sum ($3962$ scaled, $783$ in \repo{theory/revision/collision_witnesses.csv}), checked by \repo{theory/revision/check_thm513.py}. (iv)~Figure~\ref{fig:F4}: \repo{figures/gen/gen_f4_area_classes.py}, recomputed independently by \repo{review/round1-fixes/a1_fig2_kmult.py}; its pairs: \repo{review/round1-fixes/d2_explicit_pairs.py}. Figure~\ref{fig:F7}: \repo{figures/gen/gen_f7_strata.py}. (v)~The rank of $C_{27/2}$: PARI/GP, \repo{theory/diophantine/data/ranks.txt}. \texttt{code/run\_all.sh} reruns all of these except the $n=4$ search of (i), whose committed output it rechecks, and the search of (ii), which is logged.

\end{appendices}

\bibliography{references}

\end{document}
